\documentclass[11pt,a4paper]{article}

%─── Packages ─────────────────────────────────────────────────────────────────
\usepackage{amsmath,amssymb,amsthm}
\usepackage{geometry}
\usepackage{booktabs}
\usepackage[bookmarks=false]{hyperref}
\usepackage{array}
\usepackage{microtype}
\usepackage{xcolor}
\usepackage{enumitem}
\usepackage{graphicx}
\usepackage{float}
\usepackage{longtable}
\geometry{top=2.5cm,bottom=2.8cm,left=2.5cm,right=2.5cm}

%─── Theorem environments ─────────────────────────────────────────────────────
\newtheorem{theorem}{Theorem}[section]
\newtheorem{lemma}[theorem]{Lemma}
\newtheorem{proposition}[theorem]{Proposition}
\newtheorem{corollary}[theorem]{Corollary}
\theoremstyle{definition}
\newtheorem{definition}[theorem]{Definition}
\theoremstyle{remark}
\newtheorem{remark}[theorem]{Remark}
\newtheorem{conjecture}{Conjecture}

%─── Macros ───────────────────────────────────────────────────────────────────
\newcommand{\vp}{\varphi}
\newcommand{\lam}{\lambda}
\newcommand{\bst}{\beta^{*}}
\newcommand{\sinW}{\sin^{2}\!\theta_{W}}
\newcommand{\TCMB}{T_{\mathrm{CMB}}}
\newcommand{\vEW}{v_{\mathrm{EW}}}
\newcommand{\Lcond}{\Lambda_{\mathrm{cond}}}
\newcommand{\LUV}{\Lambda_{\mathrm{UV}}}
\newcommand{\MPl}{M_{\mathrm{Pl}}}
\newcommand{\LQCD}{\Lambda_{\mathrm{QCD}}}
\newcommand{\ERy}{E_{\mathrm{Ry}}}
\newcommand{\aem}{\alpha_{\mathrm{em}}}
% Tower level macro
\newcommand{\nlvl}[1]{n\!\left(#1\right)}
% IE abbreviation
\newcommand{\IE}{\mathrm{IE}}

%─── Colours ──────────────────────────────────────────────────────────────────
\definecolor{darkgreen}{RGB}{0,120,60}
\definecolor{darkred}{RGB}{160,20,20}
\definecolor{darkblue}{RGB}{0,40,140}
\definecolor{darkorange}{RGB}{180,80,0}

%══════════════════════════════════════════════════════════════════════════════
\title{\textbf{The Golden-Spiral as a Universal Mass Hierarchy Map:\\
  From Condensate Scale to Chemical Periodicity}
\\[0.6em]
{\large Companion Paper XI to the Density-Feedback
Faddeev--Niemi Hopfion Series}
\\[0.6em]
\large Preprint --- April 2026}
\author{Frederick Manfredi}

\date{2026}
%══════════════════════════════════════════════════════════════════════════════

\begin{document}
\maketitle

\paragraph*{Units and conventions.}
Natural units $\hbar = c = k_B = 1$ throughout.
Energies in eV, MeV, or GeV.
$\MPl = 1.22\times10^{19}\,\mathrm{GeV}$ (unreduced Planck mass,
consistent with Papers~I--VII).
The condensate scale is
$\Lcond = \TCMB(\pi^2/150)^{1/4} = 1.1895\times10^{-4}\,\mathrm{eV}$.
The fine-structure constant and electron mass are framework outputs:
$\aem^{-1} = 137.035998$ (Paper~IV~\cite{Paper4}, four-term formula, $5\times10^{-7}\%$) and
$m_e = 0.51100\,\mathrm{MeV}$ (Paper~IV~\cite{Paper4}).

%─── Abstract ─────────────────────────────────────────────────────────────────
\begin{abstract}
The density-feedback Faddeev--Niemi Hopfion framework
of Papers~I--X~\cite{Paper1,Paper2,Paper3,Paper4,Paper5,Paper6,Paper7,Paper8,Paper9,Paper10}
derives $\alpha_\mathrm{em}^{-1}=137.035998$ and $m_e=0.511\,\mathrm{MeV}$
from a single experimental input $\TCMB$ with no free parameters
(Paper~XII~\cite{Paper12} proves the one-parameter chain).
We show that these two derived quantities, together with the
$\vp$-tower $m_n = \Lcond\,\vp^{-2n}$, place \emph{all} known energy
scales --- particle masses, gauge-boson masses, QCD, electroweak,
UV, and Planck scales, as well as the first ionisation energies of
the chemical elements --- on a single logarithmic spiral in which
one revolution corresponds to six tower levels and the radial
coordinate is $r = \log(1 + 0.8|n|)$.

The central result is a \textbf{Corollary of Papers~III and~IV}:
the Rydberg energy $\ERy = \aem^2\,m_e/2 = 13.6057\,\mathrm{eV}$
is a \emph{fixed, calculable tower level}
$n(\ERy) = -12.102$, with no free parameters.
Since $\aem$ and $m_e$ are both derived from $\TCMB$ alone
(up to the $0.22\%$ thick-torus residual of Paper~IV),
the entire atomic energy scale is a prediction of the condensate.

Three structural observations follow from placing the
periodic table on the spiral:
(i)~the visible electromagnetic spectrum ($1.77$--$3.1\,\mathrm{eV}$)
coincides with the pure tower level
$\Lcond\,\vp^{20} = 1.799\,\mathrm{eV}$
to $1.7\%$, linking the optical window directly to the condensate
scale via the same level index as the electron mass;
(ii)~hydrogen and oxygen are near-degenerate on the spiral
($\Delta n = 0.0015$, $\Delta\IE/\IE = 0.15\%$), both anchored to
$\ERy$ --- a geometric consequence of the framework fixing the
Rydberg as a tower level;
(iii)~the Goldschmidt geochemical substitution pairs
(Mg/Ni, Fe/Ge, Fe/Co) cluster within $\Delta n < 0.004$
on the spiral, giving the first framework-level explanation for why
these pairs substitute freely in crystal lattices.
We propose the $\vp$-spiral as a supplemental organisational
principle for the periodic table, ordered by binding energy rather
than by electron configuration, and identify the testable prediction
that spiral proximity should correlate with Goldschmidt compatibility.
Spiral proximity predicts Goldschmidt substitution at AUC~$=0.857$
(Theorem~\ref{P11:thm:auc}), outperforming ionic-radius proximity alone.

Section~\ref{P11:sec:successive_IE} extends the spiral to successive
ionisation energies.
The master formula
$\nlvl{\IE_k}=\nlvl{\ERy}-\log_\vp Z_{\rm eff}(k)+\log_\vp n_{\rm qn}(k)$
(Proposition~\ref{P11:prop:master_formula}) places every IE$_k$ on the
spiral from the same two framework outputs with no additional parameters.
Two placements are examined against their chance rates: the
near-integer positions of shell-completion IEs
(Proposition~\ref{P11:prop:shell_completion}), which a census of the
tabulated ionisation energies finds to occur at or below the rate
expected for levels uniformly distributed modulo one; and the
inner-shell pair O~IE$_6$--Na~IE$_5$ at $|\Delta n|=0.0021$
(Proposition~\ref{P11:prop:inner_degen}), numerically comparable to
the H--O outer-shell result but likewise consistent with chance under
its own trial count.

The H--O near-degeneracy is examined in detail
(Proposition~\ref{P11:prop:R_equals_1} and Remark~\ref{P11:rem:oalpha_excess}).
What the framework fixes with no free parameters is the Rydberg level
itself, where hydrogen sits by definition; that oxygen lands within
$\Delta n = 0.0015$ of it is a fact of its $1s^22s^22p^4$ structure,
which analytic atomic theory (the exact hydrogenic Slater--Condon
integrals and Condon--Shortley angular algebra) reproduces to a few
per cent but not to the observed precision.  The residual $O(\aem)$
excess $\IE(\mathrm{O})/\ERy - 1 = 9.04\times10^{-4}$ is coincidentally
close to $3\aem/(8\pi)$, but that form is excluded as a mechanism by
the hydrogen ionisation energy and is the ordinary $2p^4$ correlation
energy.
\end{abstract}

\footnotesize
\noindent\textbf{Keywords:} golden ratio spiral; periodic table; ionisation energy; Rydberg energy; phi-tower; Goldschmidt compatibility; geochemistry; crystal substitution; oxygen effective nuclear charge; Slater shielding; Condon-Shortley exchange; WZW radiative correction; hydrogen-oxygen near-degeneracy; AUC; ROC analysis.
\normalsize

\tableofcontents
\bigskip

\newpage
%══════════════════════════════════════════════════════════════════════════════
\section{Introduction}
\label{P11:sec:intro}
%══════════════════════════════════════════════════════════════════════════════

The companion Papers~I--XII established a single-input framework deriving
Standard-Model parameters, QCD, gravity, and cosmology from the
condensate scale $\Lcond$ (set by $\TCMB$); the electron mass $m_e$
is derived from $\TCMB$ alone via the one-parameter chain of
Paper~XII~\cite{Paper12}.
The key structural object is the $\vp$-tower
\begin{equation}
  m_n \;=\; \Lcond\,\vp^{-2n}, \qquad n \in \mathbb{R}, \qquad
  \vp = \tfrac{1+\sqrt{5}}{2},
  \label{P11:eq:tower}
\end{equation}
which spaces successive energy scales by the golden ratio squared
$\vp^2 \approx 2.618$ per level.  The Hopf charge $Q=10$ identifies
the electron at level $|n_e| \approx 23$; the WZW T-matrix phases
place quarks and leptons; the one-loop Seeley--DeWitt coefficient
places $\MPl$ at level $|n_{Pl}| \approx 76.6$; the
UV cutoff sits at $|n_{UV}| \approx 73.6$.

All of this lives in particle physics and gravity.  The question
addressed in this paper is whether the same tower also organises
atomic and chemical physics — and whether the framework makes
testable predictions about the periodic table.

The answer is yes, in two distinct senses:
\begin{enumerate}[noitemsep]
\item \emph{Provably}, because $\aem$ and $m_e$ are framework outputs,
  the Rydberg energy $\ERy = \aem^2 m_e/2$ is a fixed tower level
  (Section~\ref{P11:sec:rydberg}).  This is a corollary of Papers~III
  and~IV with no additional assumptions.
\item \emph{Structurally}, placing the first ionisation energies
  of the elements on the tower reveals near-degeneracies
  (Section~\ref{P11:sec:near_deg}) and Goldschmidt substitution clusters
  (Section~\ref{P11:sec:goldschmidt}) that have direct chemical
  consequences.  These are observations from the framework's fixed
  outputs, not new predictions requiring new parameters.
\end{enumerate}

\paragraph{Organisation.}
Section~\ref{P11:sec:spiral} defines the $\vp$-spiral precisely.
Section~\ref{P11:sec:rydberg} proves the Rydberg corollary and places
the EM spectrum on the spiral.
Section~\ref{P11:sec:elements} maps the periodic table.
Section~\ref{P11:sec:near_deg} analyses the near-degeneracies.
Section~\ref{P11:sec:goldschmidt} connects to Goldschmidt substitution.
Section~\ref{P11:sec:successive_IE} extends the spiral to successive
ionisation energies and examines the inner-shell O--Na near-degeneracy.
Section~\ref{P11:sec:open} states open problems.
Section~\ref{P11:sec:summary} summarises.

%══════════════════════════════════════════════════════════════════════════════
\section{The $\vp$-Spiral}
\label{P11:sec:spiral}
%══════════════════════════════════════════════════════════════════════════════

\subsection{Definition}

\begin{definition}[$\vp$-tower level]
\label{P11:def:tower_level}
For any energy scale $E > 0$, its \emph{$\vp$-tower level} is
\begin{equation}
  \nlvl{E} \;=\; -\frac{\ln(E/\Lcond)}{2\ln\vp}.
  \label{P11:eq:tower_level}
\end{equation}
Equivalently, $E = \Lcond\,\vp^{-2\,n(E)}$.
Level $n=0$ is the condensate ground state $\Lcond$;
increasing $|n|$ corresponds to higher energy.
\end{definition}

The tower level is a \emph{continuous} quantity: particle masses
need not land on integers, and in general they do not (the
WZW spoke corrections shift them off pure integer levels).

\begin{remark}[Quantum origin of the tower spacing]
\label{P11:rem:quantum_tower_origin}
The spacing $\vp^2$ between successive tower levels is not merely a
display choice: Paper~X~\cite{Paper10} (Theorem~6.3 therein) proves that
the quantised condensate Hamiltonian has spectrum $E_n = \vp^{2n}E_0$
and that the $q$-deformation parameter of the resulting Macfarlane--Biedenharn
oscillator algebra is exactly $q = \vp^2$.
The tower of Table~\ref{P11:tab:levels} is therefore the semiclassical limit
of a quantum spectrum, with $\vp^2$ being the quantum of the condensate ladder.
\end{remark}

\begin{definition}[$\vp$-spiral]
\label{P11:def:spiral}
The \emph{$\vp$-spiral} is the curve
$(\theta(n), r(n))$ in cylindrical coordinates defined by
\begin{equation}
  \theta(n) = \frac{2\pi\,|n|}{N_{\rm arms}},
  \qquad
  r(n) = R_0\,\log(1 + 0.8\,|n|),
  \qquad
  y(n) = |n|\cdot\delta y,
  \label{P11:eq:spiral}
\end{equation}
where $N_{\rm arms} = 6$ is the number of tower levels per revolution,
$R_0$ is a display scale factor, and $\delta y$ is a vertical rise per
level.
The logarithmic radial coordinate $r = R_0\log(1+0.8|n|)$ compresses
the 77-level range from $\Lcond$ to $\MPl$ into a finite viewing
region while preserving the relative angular positions that encode
the physics.
\end{definition}

\begin{remark}[Choice of $N_{\rm arms}$]
\label{P11:rem:Narms_choice}
The value $N_{\rm arms} = 6$ is a display choice, not a physical one.
Any value gives the same physics; the angular groupings of elements
and particles change as $N_{\rm arms}$ varies, making the spiral an
interactive diagnostic tool.  The interactive Three.js visualisation
(Appendix~\ref{P11:app:code}) includes a live slider for $N_{\rm arms}$.
\end{remark}

\subsection{Key scales on the spiral}

Table~\ref{P11:tab:levels} lists the tower levels of the framework's
principal energy scales.

\begin{table}[h]
\centering
\caption{$\vp$-tower levels $n(E) = -\ln(E/\Lcond)/(2\ln\vp)$
  for key framework scales.
  $\Lcond = 1.1895\times10^{-4}\,\mathrm{eV}$ (single experimental input
  $\TCMB$; zero free parameters).}
\label{P11:tab:levels}
\begin{tabular}{llrr}
\toprule
Scale & Value & $n$ & Source \\
\midrule
$\Lcond$ (hub)         & $1.190\times10^{-4}\,\mathrm{eV}$ & $0.00$   & $\TCMB$ input \\
$\TCMB$                & $2.349\times10^{-4}\,\mathrm{eV}$ & $-0.71$  & Paper~I~\cite{Paper1} \\
$m_\nu$ (pred.)        & $36\,\mathrm{meV}$                & $-5.94$  & Paper~VII~\cite{Paper7} \\
$m_e$ (input)          & $0.511\,\mathrm{MeV}$             & $-23.05$ & Paper~IV~\cite{Paper4} \\
$\LQCD$                & $208\,\mathrm{MeV}$               & $-29.29$ & Paper~V~\cite{Paper5} \\
$m_\tau$               & $1777\,\mathrm{MeV}$              & $-31.52$ & Paper~I \\
$m_t$                  & $172.76\,\mathrm{GeV}$            & $-36.28$ & Paper~V \\
$\vEW$                 & $246.2\,\mathrm{GeV}$             & $-36.64$ & Paper~VI~\cite{Paper6} \\
$\LUV$                 & $6.8\times10^{17}\,\mathrm{GeV}$  & $-73.59$ & Paper~VII \\
$\MPl$                 & $1.22\times10^{19}\,\mathrm{GeV}$ & $-76.59$ & Paper~VII \\
\midrule
$\ERy = \aem^2 m_e/2$  & $13.606\,\mathrm{eV}$             & $-12.102$ & \textbf{Cor.~\ref{P11:cor:rydberg}} \\
$\Lcond\,\vp^{20}$ (Proposition~\ref{P11:prop:visible})     & $1.799\,\mathrm{eV}$              & $-10.00$  & Section~\ref{P11:sec:visible} \\
\bottomrule
\end{tabular}
\end{table}

%══════════════════════════════════════════════════════════════════════════════
\section{The Rydberg Energy as a Fixed Tower Level}
\label{P11:sec:rydberg}
%══════════════════════════════════════════════════════════════════════════════

\subsection{The Rydberg corollary}

\begin{corollary}[Rydberg as derived tower level]
\label{P11:cor:rydberg}
Within the density-feedback Hopfion framework,
the Rydberg energy $\ERy = \aem^2 m_e/2$ is a
\emph{calculable, parameter-free} tower level:
\begin{equation}
  \boxed{\nlvl{\ERy}
  \;=\; n(m_e) + \frac{\ln(2/\aem^2)}{2\ln\vp}
  \;=\; -12.102},
  \label{P11:eq:nRy}
\end{equation}
and the Rydberg energy itself is
\begin{equation}
  \boxed{
    \ERy
    \;=\; \frac{\aem^2\,m_e}{2}
    \;=\; 13.606\,\mathrm{eV}
    \;=\; \Lcond\,\vp^{24.204},
  }
  \label{P11:eq:Ry_value}
\end{equation}
where every factor is a framework output:
$\aem^{-1} = 360/\vp^2 - 3/(2\pi) + 1/(9\vp^6) - 1/(36\pi\vp^6)= 137.036$
(Paper~IV, Theorem~5.3 and Corollary~8.11, four-term, $5\times10^{-7}\%$)) and
$m_e = \Lcond\,\vp^{20}\,e^{(1600\pi-3)/400}\cdot(1+O(R_0^{-2}))$
(Paper~IV, Verlinde cancellation~\cite{Paper4}).
\end{corollary}

\begin{proof}
Direct substitution.  The tower level of $\ERy = \aem^2 m_e/2$ is
\begin{align}
  \nlvl{\ERy}
  &= -\frac{\ln(\aem^2 m_e / 2\,\Lcond)}{2\ln\vp}
   = -\frac{\ln(m_e/\Lcond)}{2\ln\vp}
     -\frac{\ln(\aem^2/2)}{2\ln\vp} \notag\\
  &= n(m_e) + \frac{\ln(2/\aem^2)}{2\ln\vp}.
  \label{P11:eq:nRy_proof}
\end{align}
Using $n(m_e) = -23.047$ (from $m_e = 0.51100\,\mathrm{MeV}$
and $\Lcond = 1.1895\times10^{-4}\,\mathrm{eV}$) and
$\aem^{-1} = 137.036$:
\begin{equation}
  \frac{\ln(2/\aem^2)}{2\ln\vp}
  = \frac{\ln(2\times137.036^2)}{2\ln\vp}
  = \frac{\ln(37558)}{2\times0.4812}
  = \frac{10.534}{0.9624}
  = 10.945,
\end{equation}
giving $\nlvl{\ERy} = -23.047 + 10.945 = -12.102$.
Numerically: $\ERy = (1/137.036)^2\times 0.51100\,\mathrm{MeV}/2
= 13.606\,\mathrm{eV}$, consistent with the measured hydrogen
ionisation energy $13.598\,\mathrm{eV}$ to $0.06\%$.
Since $\aem$ is accurate to $5\times10^{-7}\%$ (Paper~IV~\cite{Paper4}) and $m_e$
is an experimental input, the Rydberg is predicted to $0.06\%$.
\qed
\end{proof}

\begin{remark}[Physical interpretation]
\label{P11:rem:rydberg_physical}
The Rydberg is the natural unit of atomic energy — the binding energy
of the hydrogen 1s electron in the Bohr model,
$\ERy = m_e \aem^2/2 = 1\,\mathrm{Ry}$.
Corollary~\ref{P11:cor:rydberg} says that this unit does not float freely:
it is pinned to the $\vp$-tower at level $-12.102$ by the same
condensate that fixes $\TCMB$ and $m_e$.
The Bohr radius $a_0 = 1/(\aem m_e)$ and all atomic length scales
follow, pinned to the same condensate.
Chemistry, in this sense, is \emph{downstream} of the condensate.

Since Paper~X~\cite{Paper10} establishes the full canonical quantisation
of the condensate (Hilbert space $\mathcal{H}=L^2(\mathcal{C}_Q^\mathrm{red},d\mu)$,
Born rule, and spectral gap), this statement holds at the quantum level:
$\aem$ and $m_e$ are outputs of a well-defined quantum field theory,
and the Rydberg is a calculable tower level of that quantum theory.
\end{remark}

\begin{remark}[The level gap $\Delta n = 10.945$]
The separation $\Delta n = n(\ERy) - n(m_e) = 10.945$
between the electron mass and the Rydberg satisfies
\begin{equation}
  \vp^{2\Delta n}
  = \vp^{21.89}
  \approx 37\,558
  = \frac{2}{\aem^2}.
  \label{P11:eq:gap}
\end{equation}
This is not an independent identity — it is just the definition
$\Delta n = \ln(2/\aem^2)/(2\ln\vp)$ restated.
It says the factor $2/\aem^2 \approx 37\,558$ that separates
the electron rest mass from atomic binding energies
is expressed as a $\vp$-power $\vp^{21.89}$ on the spiral.
Since $\aem$ is a framework output, this ratio is predicted,
not fitted.
\end{remark}

\subsection{The visible light band as a tower level}
\label{P11:sec:visible}

\begin{proposition}[Red-light threshold as pure tower level]
\label{P11:prop:visible}
The red edge of the visible electromagnetic spectrum
($\approx 1.77\,\mathrm{eV}$) coincides with the
pure tower level at $|n|=20$:
\begin{equation}
  \Lcond\,\vp^{20}
  \;=\; 1.1895\times10^{-4}\,\mathrm{eV}\times\vp^{20}
  \;=\; 1.799\,\mathrm{eV},
  \label{P11:eq:red_light}
\end{equation}
accurate to $1.7\%$.
\end{proposition}

\begin{proof}
$\vp^{20} = 15127.000$ (exact Lucas number $L_{20}$);
$\Lcond\times 15127.000 = 1.799\,\mathrm{eV}$.
The measured red-light threshold is $\approx 1.77\,\mathrm{eV}$.
Deviation: $|(1.799-1.770)/1.770| = 1.6\%$.
\qed
\end{proof}

\begin{remark}[Two uses of level 20]
\label{P11:rem:level20_dual_use}
The same level index $n=20$ appears in two different formulas:
\begin{align}
  \text{Red-light threshold:} &\quad \Lcond\,\vp^{20} = 1.799\,\mathrm{eV}, \\
  \text{Electron mass:} &\quad m_e = \Lcond\,\vp^{20}\,e^{4\pi-3/400}\cdot(1+O(R_0^{-2})),
\end{align}
where $e^{4\pi-3/400} \approx 2.87\times10^5$.
The pure tower level at $n=20$ is the red-light threshold.
The physical electron mass sits at the same tower index but is lifted
by the WZW Chern--Simons exponential $e^{4\pi}$ from this threshold
by a factor of $\approx 286\,000$.
In other words: the optical window ``knows about'' the electron mass
tower level, but not its WZW lift.
\end{remark}

\subsection{The full EM spectrum on the spiral}

Table~\ref{P11:tab:em} lists the tower levels of the electromagnetic
spectrum.  The chemistry band (first ionisation energies,
$\approx 4$--$25\,\mathrm{eV}$, levels $-10.8$ to $-12.7$) sits
immediately above the visible band on the spiral, anchored by the
Rydberg at $n=-12.10$.  The co-location of visible light and
chemistry on the spiral is a geometric expression of the
physical fact that visible photons couple resonantly to atomic
electrons: both are anchored to $\ERy$.

\begin{table}[h]
\centering
\caption{EM spectrum and chemistry band on the $\vp$-spiral.}
\label{P11:tab:em}
\begin{tabular}{llr}
\toprule
Band / feature & Energy & Tower level $n$ \\
\midrule
Radio (AM, $540\,\mathrm{kHz}$) & $2.2\times10^{-9}\,\mathrm{eV}$ & $+5.6$ \\
Microwave peak ($1\,\mathrm{mm}$) & $1.2\times10^{-3}\,\mathrm{eV}$ & $-2.2$ \\
Far infrared                     & $12\,\mathrm{meV}$               & $-4.8$ \\
Near infrared edge               & $0.8\,\mathrm{eV}$               & $-9.2$ \\
\textbf{Red light} (threshold)   & $\mathbf{1.77\,\mathrm{eV}}$     & $\mathbf{-9.98}$ \\
$\Lcond\,\vp^{20}$ (pure tower)  & $1.799\,\mathrm{eV}$             & $-10.00$ \\
\textbf{Blue light}              & $\mathbf{3.1\,\mathrm{eV}}$      & $\mathbf{-10.57}$ \\
Ultraviolet                      & $6.2\,\mathrm{eV}$               & $-11.29$ \\
\textbf{Rydberg / H IE}          & $\mathbf{13.606\,\mathrm{eV}}$   & $\mathbf{-12.102}$ \\
Soft X-ray                       & $100\,\mathrm{eV}$               & $-14.18$ \\
Hard X-ray                       & $10\,\mathrm{keV}$               & $-18.96$ \\
Gamma                            & $1\,\mathrm{MeV}$                & $-23.74$ \\
\midrule
Chemistry band (IE range)        & $4$--$25\,\mathrm{eV}$           & $-10.8$ to $-12.7$ \\
\bottomrule
\end{tabular}
\end{table}

%══════════════════════════════════════════════════════════════════════════════
\section{Chemical Elements on the $\vp$-Spiral}
\label{P11:sec:elements}
%══════════════════════════════════════════════════════════════════════════════

\subsection{Placement by first ionisation energy}

We place each element on the spiral at the tower level
$n(\IE_Z) = -\ln(\IE_Z/\Lcond)/(2\ln\vp)$ of its
first ionisation energy $\IE_Z$.
The first IE is the energy cost of removing the outermost
electron — it is set by $\aem$, $m_e$, and the effective
nuclear charge $Z_\mathrm{eff}$ that the outermost electron
experiences after screening by inner-shell electrons.
Since $\aem$ and $m_e$ are framework outputs
(Section~\ref{P11:sec:rydberg}), the tower level of $\IE_Z$ is
determined by $Z_\mathrm{eff}$ alone.

\begin{remark}[What the spiral measures]
\label{P11:rem:spiral_measures}
Two elements at the same tower level have outermost electrons
bound with the same energy.  This is a statement about their
effective nuclear charge, independent of their position in
the conventional periodic table.  The spiral thus provides a
binding-energy organisation of chemistry, complementary to
the electron-configuration organisation of the standard table.
\end{remark}

\subsection{Group structure on the spiral}
\label{P11:sec:group_structure}

Table~\ref{P11:tab:elements} lists the tower levels of selected elements,
grouped by periodic table group.
First ionisation energies are from the NIST Atomic Spectra
Database~\cite{NIST_ASD}; the effective nuclear charges
are derived from these via the hydrogenic formula
(see Slater~\cite{Slater1930} for the screening approximation).
Measured particle masses are from the PDG~\cite{PDG2024}.

\begin{longtable}{llrrp{5cm}}
\caption{Selected elements on the $\vp$-spiral,
  placed by first ionisation energy.
  Source: NIST Atomic Spectra Database.}
\label{P11:tab:elements}\\
\toprule
Element & Group & IE (eV) & $n$ & Notes \\
\midrule
\endfirsthead
\toprule
Element & Group & IE (eV) & $n$ & Notes \\
\midrule
\endhead
\bottomrule
\endfoot
Cs  &  1 (alkali)  &  3.894 & $-10.80$ & Lightest alkali IE \\
Fr  &  1           &  4.073 & $-10.85$ & \\
Rb  &  1           &  4.177 & $-10.87$ & \\
K   &  1           &  4.341 & $-10.91$ & \\
Na  &  1           &  5.139 & $-11.09$ & \\
Li  &  1           &  5.392 & $-11.14$ & \\
\midrule
Ba  &  2 (alk. earth) & 5.212 & $-11.11$ & \\
Ca  &  2              & 6.113 & $-11.27$ & \\
Sr  &  2              & 5.695 & $-11.20$ & \\
Mg  &  2              & 7.646 & $-11.50$ & Near-degenerate with Ni \\
Be  &  2              & 9.323 & $-11.71$ & \\
\midrule
Sc  &  3 (d-block)    & 6.561 & $-11.34$ & \\
Mn  &  7              & 7.434 & $-11.47$ & \\
Fe  &  8              & 7.902 & $-11.54$ & Near-degenerate with Ge, Co \\
Co  &  9              & 7.881 & $-11.53$ & Near-degenerate with Fe \\
Ni  & 10              & 7.640 & $-11.50$ & Near-degenerate with Mg \\
Cu  & 11              & 7.727 & $-11.51$ & \\
Zn  & 12              & 9.394 & $-11.72$ & \\
Au  & 11 (period 6)   & 9.226 & $-11.70$ & Near-degenerate with Zn \\
\midrule
Ge  & 14 (p-block)    & 7.900 & $-11.54$ & Near-degenerate with Fe \\
Si  & 14              & 8.152 & $-11.57$ & \\
Sn  & 14              & 7.344 & $-11.46$ & \\
C   & 14              & 11.260 & $-11.91$ & \\
\midrule
N   & 15              & 14.534 & $-12.17$ & \\
P   & 15              & 10.487 & $-11.83$ & \\
O   & 16              & 13.618 & $-12.10$ & \textbf{Near-degenerate with H} \\
S   & 16              & 10.360 & $-11.82$ & \\
\midrule
F   & 17 (halogen)    & 17.423 & $-12.36$ & \\
Cl  & 17              & 12.968 & $-12.05$ & \\
Br  & 17              & 11.814 & $-11.96$ & \\
I   & 17              & 10.451 & $-11.83$ & \\
\midrule
He  & 18 (noble)      & 24.587 & $-12.72$ & Highest IE \\
Ne  & 18              & 21.565 & $-12.58$ & \\
Ar  & 18              & 15.760 & $-12.25$ & \\
Kr  & 18              & 13.999 & $-12.13$ & \\
Xe  & 18              & 12.130 & $-11.98$ & \\
Rn  & 18              & 10.748 & $-11.86$ & \\
\midrule
\textbf{H}   & 1      & \textbf{13.598} & $\mathbf{-12.101}$ & Rydberg exactly \\
\end{longtable}

The groups visible in the conventional periodic table
are partially preserved as angular arcs on the spiral
(alkalis cluster near $n \approx -10.8$ to $-11.1$;
noble gases span $-11.86$ to $-12.72$),
but elements from different groups also cluster tightly
when their effective nuclear charges coincide.

%══════════════════════════════════════════════════════════════════════════════
\section{Near-Degeneracies and Their Chemical Consequences}
\label{P11:sec:near_deg}
%══════════════════════════════════════════════════════════════════════════════

\subsection{Statement of the near-degeneracies}

Table~\ref{P11:tab:degen} lists the four tightest near-degeneracies
found in the 118-element dataset.

\begin{table}[h]
\centering
\caption{Near-degeneracies on the $\vp$-spiral.
  $\Delta n = |n(E_1) - n(E_2)|$.}
\label{P11:tab:degen}
\begin{tabular}{llrrrrp{4.5cm}}
\toprule
Pair & Groups & IE$_1$ (eV) & IE$_2$ (eV) & $\Delta$IE (eV)
  & $\Delta n$ & Chemical consequence \\
\midrule
\textbf{Fe--Ge} & 8, 14 & 7.902 & 7.900 & 0.002 & 0.00026
  & Ge semiconductor doped with Fe impurities \\
\textbf{Mg--Ni} & 2, 10 & 7.646 & 7.640 & 0.006 & 0.00082
  & Olivine $(Mg,Ni)_2$SiO$_4$ crystal substitution \\
\textbf{H--O} & 1, 16 & 13.598 & 13.618 & 0.020 & 0.00153
  & Water formation; both at the Rydberg \\
\textbf{Fe--Co} & 8, 9 & 7.902 & 7.881 & 0.021 & 0.00276
  & Fe--Co alloys; Heusler compounds \\
\bottomrule
\end{tabular}
\end{table}

\subsection{The H--O near-degeneracy}
\label{P11:subsec:ho}

Hydrogen (IE $= 13.598\,\mathrm{eV}$) and oxygen (IE $= 13.618\,\mathrm{eV}$)
are separated by $\Delta n = 0.00153$ on the spiral
--- a relative energy difference of $0.15\%$.
Both sit at the Rydberg level $n \approx -12.10$.

\paragraph{Why hydrogen is there.}
By definition: the hydrogen 1s binding energy is the Rydberg,
$\IE_\mathrm{H} = \aem^2 m_e/2 = \ERy = 13.606\,\mathrm{eV}$.
(The measured value $13.598\,\mathrm{eV}$ includes the
proton recoil correction $m_e/(m_e + m_p) \approx 0.9995$.)

\paragraph{Why oxygen is there.}
Oxygen's outermost electron is a 2p electron with principal
quantum number $n_\mathrm{qn}=2$ and effective nuclear charge
defined by the hydrogenic formula $\IE = Z_\mathrm{eff}^2\,\ERy/n_\mathrm{qn}^2$:
\begin{equation}
  Z_\mathrm{eff}(\mathrm{O})
  = n_\mathrm{qn}\sqrt{\IE_\mathrm{O}/\ERy}
  = 2\sqrt{13.618/13.606}
  = 2.001.
  \label{P11:eq:Zeff_O}
\end{equation}
So oxygen's effective nuclear charge is $Z_\mathrm{eff} \approx 2.001$
in units of the Rydberg scale --- essentially exactly $n_\mathrm{qn}$.
The condition $Z_\mathrm{eff}=n_\mathrm{qn}$ is precisely the
statement $\IE(\mathrm{O}) = \ERy$: the many-body screening of O's
$1s^22s^22p^4$ configuration reduces the apparent nuclear charge
until the outermost electron is bound with the same energy as
hydrogen's 1s electron.
A Slater-rule estimate~\cite{Slater1930} for the $2p$ electron gives
shielding $\sigma = 0.85\times2\;(\text{1s}) + 0.35\times5\;(\text{other }n{=}2) = 3.45$
(the $2s$ and $2p$ electrons form a single Slater group, screening one
another at $0.35$), yielding
$Z_\mathrm{eff}^\mathrm{Slater} = 8-3.45 = 4.55$.
Inserted into the one-electron orbital formula this gives
$\IE = 4.55^2\,\ERy/4 = 70.4\,\mathrm{eV}$, five times too large ---
but that formula is not Slater's prescription.  Used as intended, as a
difference of total energies
$E=-\ERy\sum_i (Z_{\mathrm{eff},i}/n^*_i)^2$ between O and O$^+$, the
same rules give $\IE = 14.17\,\mathrm{eV}$, within $4\%$ of
experiment ($|\Delta n| = 0.041$).
The value $Z_\mathrm{eff} = 2.001$ at the precision required here
still comes only from a full ab initio treatment
(Section~\ref{P11:sec:zeff_derivation}).
The small discrepancy from 2.000 reflects beyond-Slater correlation
effects.

\paragraph{Framework interpretation.}
The framework fixes $\aem$ and $m_e$, hence fixes $\ERy = 13.606\,\mathrm{eV}$
at tower level $-12.102$ with no free parameters.
The coincidence of H and O at this level is thus:
H is there by definition;
O is there because $Z_\mathrm{eff}(\mathrm{O}) \approx 2$
(Section~\ref{P11:sec:zeff_derivation} shows this follows from the condensate
via the 8-electron Schrödinger equation with framework-fixed $\aem$).
What the framework contributes is that the Rydberg level is
\emph{fixed on the tower} --- chemistry cannot shift it
without changing $\aem$ or $m_e$.

\subsection{The oxygen ionisation energy and the condensate}
\label{P11:sec:zeff_derivation}

We trace how the oxygen ionisation energy, and hence its placement
at the Rydberg level, follows from the condensate's output values
$\aem$ and $m_e$ through the parameter-free eight-electron atom ---
to the few-per-cent accuracy of an analytic treatment, the residual
precision being ordinary many-body structure.

\paragraph{Step 1: Condensate fixes all atomic input parameters.}
The condensate delivers two quantities relevant to atomic physics
(Papers~III and~IV):
\begin{align}
  \aem^{-1} &= 137.035\,998 \qquad (\text{Paper~IV, four-term, }5\times10^{-7}\%), \\
  m_e       &= 0.51100\,\mathrm{MeV} \qquad (\text{Paper~IV input}).
\end{align}
These fix the Bohr radius $a_0 = 1/(\aem m_e)$, the Rydberg $\ERy = \aem^2 m_e/2$,
and therefore the complete Coulomb Hamiltonian for any $Z$-electron atom:
\begin{equation}
  H_Z = \sum_{i=1}^Z \left(-\frac{\nabla_i^2}{2m_e} - \frac{Z\aem}{r_i}\right)
       + \sum_{i<j} \frac{\aem}{r_{ij}}.
  \label{P11:eq:HZ}
\end{equation}
Once $\aem$ and $m_e$ are fixed, $H_8$ (the oxygen Hamiltonian) has
no remaining free parameters.

\paragraph{Step 2: The condition $Z_\mathrm{eff}(\mathrm{O})=2$.}
The effective nuclear charge $Z_\mathrm{eff}$ for the outermost ($2p$) electron
of oxygen is defined operationally by the hydrogenic inversion of the
first ionisation energy:
\begin{equation}
  Z_\mathrm{eff}(\mathrm{O}) \;\equiv\; n_\mathrm{qn}\sqrt{\IE(\mathrm{O})/\ERy}
  \;=\; 2\sqrt{13.618/13.606} \;=\; 2.001.
  \label{P11:eq:Zeff_def}
\end{equation}
The condition $Z_\mathrm{eff}=n_\mathrm{qn}=2$ is precisely $\IE(\mathrm{O})=\ERy$,
i.e.\ oxygen's first ionisation energy equals the hydrogen ionisation energy.

\paragraph{Step 3: Why $\IE(\mathrm{O}) \approx \ERy$.}
The Hartree--Fock ground-state energy of oxygen's $1s^22s^22p^4$ configuration
can be written as
\begin{equation}
  E_0(\mathrm{O}) = \sum_i \varepsilon_i - \tfrac{1}{2}\sum_{i\neq j}J_{ij}
  + \tfrac{1}{2}\sum_{i\neq j}K_{ij},
\end{equation}
where $\varepsilon_i$ are orbital energies, $J_{ij}$ Coulomb integrals,
and $K_{ij}$ exchange integrals.  Koopmans' theorem gives
$\IE(\mathrm{O}) \approx -\varepsilon_{2p}$.

The orbital energy $\varepsilon_{2p}$ depends on $\aem$, $m_e$, and
the self-consistent screening by the remaining 7 electrons.
Analytic orbital treatments bracket the measurement rather than
reaching it: Koopmans' theorem (the frozen-orbital HF eigenvalue)
overestimates $\IE(\mathrm{O})$ by about $26\%$, the Hartree--Fock
total-energy difference underestimates it by about $13\%$, and the gap
between them is correlation energy.  The result lands near $\ERy$ to a
few per cent, not to $0.1\%$.

The physical mechanism is a cancellation between competing many-body effects:
\begin{enumerate}[noitemsep]
\item \textbf{Increased nuclear attraction:} $Z=8$ pulls the 2p electron
  inward more than hydrogen's $Z=1$.
\item \textbf{Inner-shell screening:} the $1s^2 2s^2$ core shields $\approx 3.5$
  units of charge, leaving effective $Z_\mathrm{eff}^\mathrm{core} \approx 4.5$.
\item \textbf{Intra-$2p$ pairing penalty:} in $2p^4$ (beyond the
  half-filled $2p^3$ of nitrogen) the fourth electron must doubly
  occupy an orbital, at the cost of the pair repulsion that lowers
  $\IE(\mathrm{O})$ below $\IE(\mathrm{N})$ despite the larger nuclear
  charge.  This penalty is $\approx 0.86\,\mathrm{eV}$
  (Remark~\ref{P11:rem:pairing_penalty}), reproducing the measured
  nitrogen--oxygen anomaly to $6\%$.
\end{enumerate}
The net effect of (1)--(3) places $\IE(\mathrm{O})$ within a few per
cent of $\ERy$: inner-shell screening brings the bare $Z=8$ attraction
down to an orbital effective charge $\approx4.5$, and the pairing
penalty of (3) supplies the further drop onto the Rydberg scale.  The
agreement is a genuine feature of the $1s^22s^22p^4$ configuration, not
an exact cancellation; the residual precision of the H--O coincidence
($0.15\%$) is beyond what this analytic accounting reaches.

\paragraph{Step 4: The condensate origin.}
Since $H_8$ in~\eqref{P11:eq:HZ} has no free parameters once $\aem$ and
$m_e$ are fixed, and since the condensate (Papers~III and~IV)
determines both, the approximate equality $\IE(\mathrm{O})\approx\ERy$
--- and hence oxygen's placement at the Rydberg level --- follows from
the framework outputs through the parameter-free eight-electron
problem:
\[
  \TCMB, m_e
  \;\xrightarrow{\text{Papers I--X}}\;
  \aem,\, m_e
  \;\xrightarrow{\text{Eq.~\eqref{P11:eq:HZ}}}\;
  H_8
  \;\xrightarrow{\text{atomic structure}}\;
  \IE(\mathrm{O}) \approx \ERy .
\]
The final step is a solved problem in quantum chemistry; full
numerical Hartree--Fock with correlation reproduces
$\IE(\mathrm{O})$~\cite{Clementi1963}, while the analytic content
below fixes it to a few per cent.  The statement $Z_\mathrm{eff}=2$ is
this approximate equality re-expressed
(Eq.~\eqref{P11:eq:Zeff_O}), not a sharper independent result.

The parameter-free content that does not require numerical HF is the
angular algebra of the configuration, which fixes the structure of
$\IE(\mathrm{O})$ to the accuracy stated below.

\paragraph{Step 5: Level-1 --- $\IE(\mathrm{O})\approx\ERy$ from the
$\aem$-independent eight-electron problem.}

The dimensionless ratio $\IE(\mathrm{O})/\ERy$ is a pure number determined
entirely by $Z=8$ and the electron configuration $1s^22s^22p^4$,
\emph{independent of} $\aem$.
To see this: in atomic units $\hbar=m_e=a_0=1$, all distances are
measured in Bohr radii $a_0 = 1/(\aem m_e)$ and all energies in
$2\ERy = \aem^2 m_e$.
The many-body Hamiltonian~\eqref{P11:eq:HZ} becomes
\begin{equation}
  \frac{H_Z}{2\ERy}
  \;=\; \sum_{i=1}^Z
        \!\left(-\tfrac{1}{2}\nabla_i^2 - \frac{Z}{r_i}\right)
  \;+\; \sum_{i<j}\frac{1}{r_{ij}},
  \label{P11:eq:HZ_au}
\end{equation}
in which $\aem$ no longer appears.
The eigenvalues of~\eqref{P11:eq:HZ_au} are pure dimensionless numbers;
multiplying by $2\ERy$ restores the physical units.
In particular,
\begin{equation}
  \frac{\IE(\mathrm{O})}{\ERy}
  \;=\; -2\,\frac{\varepsilon_{2p}}{2\ERy}
  \;\equiv\; R(\text{O}),
  \label{P11:eq:R_O}
\end{equation}
where $R(\mathrm{O})$ is a pure number determined by $Z=8$ and the
$1s^22s^22p^4$ configuration.

\begin{proposition}[Level-1: $R(\mathrm{O})\approx1$ from the eight-electron problem]
\label{P11:prop:R_equals_1}
The dimensionless ratio for the 2p ionisation of oxygen in the
$1s^22s^22p^4\,{}^3P$ ground state,
\begin{equation}
  R(\mathrm{O}) \;\equiv\; \frac{\IE(\mathrm{O})}{\ERy},
  \label{P11:eq:R_equals_1}
\end{equation}
is fixed by $Z=8$ and the configuration alone, and equals unity to the
few-per-cent accuracy of an analytic orbital treatment:
$\boxed{R(\mathrm{O}) = 1 + O(\text{few}\,\%)}$.  The measured value is
$R(\mathrm{O}) = 1.000\,90$; the residual $9.04\times10^{-4}$ excess is
$2p^4$ correlation energy --- itself condensate-downstream through the
framework-fixed Hamiltonian, but not a coupling shift beyond it
(Remark~\ref{P11:rem:oalpha_excess}).
\end{proposition}

\begin{proof}
We work in atomic units~\eqref{P11:eq:HZ_au}; ``Ry'' below means $2\ERy$.
By Koopmans' theorem, $\IE(\mathrm{O}) = -\varepsilon_{2p}$, and writing
the orbital energy in hydrogenic form defines an effective charge,
\begin{equation}
  -\varepsilon_{2p}
  \;=\; \frac{Z_\mathrm{eff}^2}{n_\mathrm{qn}^2}\,\ERy,
  \qquad
  Z_\mathrm{eff} \;\equiv\; Z - \sigma_\mathrm{tot},
  \label{P11:eq:eps_Zeff}
\end{equation}
so that $R(\mathrm{O})=1$ is exactly $Z_\mathrm{eff}=n_\mathrm{qn}=2$.
Equation~\eqref{P11:eq:eps_Zeff} is a change of variables rather than a
derivation: it \emph{defines} $Z_\mathrm{eff}$ from the energy, so
``$Z_\mathrm{eff}(\mathrm{O})=2$'' and ``$\IE(\mathrm{O})=\ERy$'' are
the same statement.  Establishing~\eqref{P11:eq:R_equals_1} requires
$\sigma_\mathrm{tot}$ from a prescription independent of the measured
energy.

\textbf{(a) Independent screening prescriptions.}
Two are standard.  Slater's rules~\cite{Slater1930} place the $2s$ and
$2p$ electrons in a single screening group, each screening the others
at $0.35$, with the $1s$ pair at $0.85$:
\begin{equation}
  \sigma_\mathrm{Slater}
  = 2\times0.85\;(1s^2) + 5\times0.35\;(\text{other } n{=}2)
  = 3.45,
  \qquad
  Z_\mathrm{eff}^\mathrm{Slater} = 4.55 .
  \label{P11:eq:sigma_slater}
\end{equation}
The self-consistent Hartree--Fock screening constants of Clementi
\emph{et al.}~\cite{Clementi1963} give, for oxygen,
$Z_\mathrm{eff}(2s)=4.4916$ and $Z_\mathrm{eff}(2p)=4.4532$.  The two
orbitals differ by $0.038$, confirming Slater's grouping directly: $2s$
and $2p$ see essentially the same effective charge, and $2s$ does not
act as an inner shell towards $2p$.  Taking the $1s$ pair at $0.85$
each, the Clementi value implies $0.369$ of screening per $n=2$
electron against Slater's $0.35$, a five per cent refinement.  The two
prescriptions agree on $Z_\mathrm{eff}(2p)$ to two per cent.

\textbf{(b) Neither prescription gives $Z_\mathrm{eff}=2$, and none can.}
Both place $Z_\mathrm{eff}(\mathrm{O},2p)$ near $4.5$.  The value
$\sigma_\mathrm{tot}=6.00$ required by~\eqref{P11:eq:eps_Zeff} lies
outside the range any screening decomposition can reach: with the $1s$
pair at $0.85$ and $2p$--$2p$ at $0.35$, each $2s$ electron would have
to screen $1.625$ units of nuclear charge, and a single electron cannot
screen more than one; granting the $1s$ pair the maximum $1.0$ each
still requires $1.475$.  Consequently $Z_\mathrm{eff}=2$ is not an
orbital screening constant but the measured ionisation energy
re-expressed through~\eqref{P11:eq:eps_Zeff}.

\textbf{(c) The parameter-free content is angular, not empirical.}
Both prescriptions above are fitted: Slater's constants $0.35$, $0.85$
and his effective quantum numbers $n^*$ are tuned to spectroscopic
data, and the Clementi values are numerical Hartree--Fock output.  The
combinatorial content of the configuration lies instead in the angular
algebra, which involves no fitted quantity and no dependence on
$\aem$.  For hydrogenic radial functions the Slater--Condon integrals
are exact rationals times $Z\ERy$; those governing the $2p$ shell are
\begin{equation}
  F^0(2p,2p) = \tfrac{93}{256}\,Z\,\ERy,
  \qquad
  F^2(2p,2p) = \tfrac{45}{256}\,Z\,\ERy,
  \label{P11:eq:slater_condon}
\end{equation}
and the angular coefficients of the $^3P$ term are fixed by Wigner
algebra, giving $K_{2p,2p} = \tfrac{3}{25}F^2(2p,2p)$.  Together with
the one-body term $-Z_\mathrm{eff}^2/(2n_\mathrm{qn}^2)$, whose
$1/n_\mathrm{qn}^2$ scaling is what makes the condition
$Z_\mathrm{eff}=n_\mathrm{qn}$ meaningful at all, these fix the
structure of $\IE(\mathrm{O})$ from $Z=8$ and the occupation
$1s^22s^22p^4$ alone.

\textbf{(d) Accuracy attainable at this order.}
What such a calculation cannot supply is precision.  For oxygen the
orbital-based analytic routes bracket the measurement rather than
reproducing it: Koopmans' theorem gives $-\varepsilon_{2p} =
17.2\,\mathrm{eV}$ ($+26\%$), the Hartree--Fock total-energy
difference gives $11.9\,\mathrm{eV}$ ($-13\%$), and the gap between
them is correlation energy, which no orbital scheme captures.  For
comparison, the fitted screening constants used as total-energy
differences,
\begin{equation}
  E = -\ERy \sum_i \Bigl(\frac{Z_{\mathrm{eff},i}}{n^*_i}\Bigr)^{\!2},
  \qquad
  \IE_\mathrm{Slater}(\mathrm{O}) = E(\mathrm{O}^+) - E(\mathrm{O})
  = 14.17\,\mathrm{eV},
  \label{P11:eq:IE_slater_total}
\end{equation}
land within four per cent, but on fitted input and so serve as a
consistency check rather than a derivation.
Equation~\eqref{P11:eq:R_equals_1} is therefore established at the
order the configuration's angular structure fixes --- some tens of per
cent --- and is verified empirically at
$\IE(\mathrm{O})/\ERy = 1.000\,90$.  Its precision is not derived
here.
\qed
\end{proof}

\begin{remark}[The $2p^4$ pairing penalty and the nitrogen--oxygen anomaly]
\label{P11:rem:pairing_penalty}
Slater's rules carry no exchange or spin dependence, so
equation~\eqref{P11:eq:IE_slater_total} cannot see the effect that
brings oxygen's first ionisation energy below nitrogen's despite the
larger nuclear charge: in $2p^4$ the fourth electron must doubly
occupy an orbital already singly occupied, at the cost of the pair
repulsion absent in the half-filled $2p^3$ shell.  That anomaly is
measured directly,
\begin{equation}
  \IE(\mathrm{N}) - \IE(\mathrm{O})
  = 14.534 - 13.618 = 0.916\,\mathrm{eV},
  \label{P11:eq:NO_anomaly}
\end{equation}
and the Condon--Shortley angular algebra~\cite{Condon1935} reproduces
it: with the angular coefficient $3/25$ of $K_{2p,2p}$ and the hydrogenic
radial integral $F^2(2p,2p) = (45/256)Z_\mathrm{eff}\ERy$
of~\eqref{P11:eq:slater_condon} evaluated at $Z_\mathrm{eff}=2$,
\begin{equation}
  \tfrac{3}{2}K_{2p,2p} = \tfrac{81}{2560}\,Z_\mathrm{eff}\,\ERy
  = 0.861\,\mathrm{eV},
  \label{P11:eq:pairing_K}
\end{equation}
within $6\%$ of~\eqref{P11:eq:NO_anomaly}.  The pairing penalty is
therefore the physical effect that lowers $\IE(\mathrm{O})$ onto the
Rydberg scale, and it is of the right size.

It does not combine additively with~\eqref{P11:eq:IE_slater_total} to
give a sharper prediction.  The fitted-screening estimate lies $0.55\,$eV above the measured
value and the pairing term is $0.86\,$eV, so subtracting one from the
other returns $13.31\,\mathrm{eV}$, now $2.3\%$ \emph{below}
experiment: the two parametrisations are not
disjoint contributions, since Slater's electrostatic constants are
fitted to configurations in which pairing is already present on
average.  What the two establish jointly is that ordinary atomic
structure brackets $\IE(\mathrm{O})$ around $\ERy$ from both sides at
the few-per-cent level --- the coincidence is a real feature of the
$1s^22s^22p^4$ configuration, not an accident of arithmetic --- while
neither route approaches the precision of the observed agreement.
\end{remark}

\begin{remark}[Precision of the H--O coincidence]
\label{P11:rem:sigma_xc_origin}
Equation~\eqref{P11:eq:IE_slater_total} accounts for
$R(\mathrm{O})=1$ to four per cent, $|\Delta n| = 0.041$.  The measured
coincidence is far tighter: $\IE(\mathrm{H})=13.598\,\mathrm{eV}$ and
$\IE(\mathrm{O})=13.618\,\mathrm{eV}$ differ by $0.15\%$,
$|\Delta n| = 0.0015$, some twenty-seven times closer than the
screening calculation predicts.  That excess precision is not derived
here, and is the substance of the near-degeneracy.

Two accounts of it are unavailable.  A decomposition
$\sigma_\mathrm{tot} = \sigma_\mathrm{Slater} + \sigma_\mathrm{xc}$
cannot supply it, because $\sigma_\mathrm{xc}$ would be defined as the
residual $\sigma_\mathrm{tot}-\sigma_\mathrm{Slater}$, making its
closure an identity rather than a check --- and, by the ceiling
argument above, no admissible $\sigma_\mathrm{xc}$ reaches
$\sigma_\mathrm{tot}=6.00$ in any case.  Nor does the $2p^4$ exchange
penalty account for the gap quantitatively: with the Condon--Shortley
angular coefficient $3/25$ and the hydrogenic
$F^2(2p,2p)=(45/256)Z_\mathrm{eff}\ERy$ at $Z_\mathrm{eff}=2$, the
exchange term is $(3/2)K_{2p,2p} = 0.86\,\mathrm{eV}$, which at
$\mathrm{d}\IE/\mathrm{d}\sigma = 2Z_\mathrm{eff}\ERy/n_\mathrm{qn}^2
= 13.6\,\mathrm{eV}$ per unit of screening corresponds to
$\Delta\sigma \approx 0.08$.

What the framework contributes is unaffected: the Rydberg energy sits
at tower level $-12.102$ with no free parameters once $\aem$ and $m_e$
are fixed (Corollary~\ref{P11:cor:rydberg}), and hydrogen sits there by
definition.  That oxygen sits within $0.0015$ of the same level is an
atomic-structure fact, approximately explained by ordinary screening
and exactly explained by neither.
\end{remark}

\paragraph{Step 6: the $O(\aem)$ excess and a candidate closed form.}

Proposition~\ref{P11:prop:R_equals_1} places $\IE(\mathrm{O})$ at the
Rydberg energy to the few-per-cent accuracy of the angular algebra.
The measured value exceeds $\ERy$ by $9.04\times10^{-4}$, an $O(\aem)$
quantity.  It is natural to ask whether the condensate fixes this
excess; the following records the candidate form and why it is not
adopted.

\begin{remark}[Why the $O(\aem)$ excess is not a condensate correction]
\label{P11:rem:oalpha_excess}
The measured first ionisation energy of oxygen exceeds the Rydberg
energy by
\begin{equation}
  \frac{\IE(\mathrm{O})}{\ERy} - 1 = 9.04\times10^{-4},
  \label{P11:eq:IE_O_excess}
\end{equation}
an $O(\aem)$ quantity, and the condensate value $\aem = 1/137.036$
makes
\begin{equation}
  \frac{k\,\aem}{8\pi}\bigg|_{k=3} = \frac{3\aem}{8\pi} = 8.71\times10^{-4}
  \label{P11:eq:IE_O_closed}
\end{equation}
a candidate closed form for it, reproducing the
excess~\eqref{P11:eq:IE_O_excess} to $3.7\%$ (a residual of $0.003\%$
of the total ionisation energy).  The coefficient has a motivating
reading: $k/(8\pi) = k/(4\pi\,\Delta Q)$ with $\Delta Q = 2$, the
Chern--Simons barrier of a two-unit topological transition being
$4\pi\,\Delta Q$ (Paper~XVI~\cite{Paper16}, where $\Delta Q = 3-1 = 2$
gives the $8\pi$ spoke of the constituent quark mass), and removing one
electron from the lepton-sector $\QH=2$ ground state
(Paper~I~\cite{Paper1}) is a $\Delta Q = 2$ event --- mirroring the
Chern--Simons factor $1/(4\pi)$ of the electron-mass formula of
Paper~IV~\cite{Paper4}.

The reading is nonetheless excluded as a physical mechanism by its own
universality.  Every single-electron ionisation is a $\Delta Q = 2$
event, so the correction would apply to hydrogen identically; but
$\IE(\mathrm{H})$ is accounted for by the reduced-mass factor alone,
$\ERy\,(1+m_e/m_p)^{-1} = 13.5983\,\mathrm{eV}$ against the measured
$13.5984\,\mathrm{eV}$ (residual $11\,\mathrm{ppm}$), and adding
$3\aem/(8\pi)\,\ERy = 11.9\,\mathrm{meV}$ would spoil this by
$860\,\mathrm{ppm}$, eighty times the residual.  A genuine QED
radiative shift at $Z=8$ is of order the Lamb scale
($\sim10^{-6}\,\mathrm{eV}$), three orders below the $11\,\mathrm{meV}$
here.  The excess~\eqref{P11:eq:IE_O_excess} is therefore many-body
structure --- the $2p^4$ correlation energy that separates Koopmans
from $\Delta$SCF (Proposition~\ref{P11:prop:R_equals_1}).  It is
condensate physics in the same parametric sense as the entire Coulomb
Hamiltonian, which is fixed once $\aem$ and $m_e$ are
(Step~1); what the hydrogen test excludes is the stronger claim that
the condensate adds a coupling shift \emph{beyond} that Hamiltonian,
universal across elements.  The agreement with $3\aem/(8\pi)$, closer
than the alternative $2\aem/(5\pi)$ ($2.8\%$) only by chance, is a
coincidence of the right order, not such a shift.  Whether the
condensate's structure instead surfaces \emph{through} the many-body
correlation dynamics --- which vanish for hydrogen and so evade this
test --- is a separate and open question, of the same kind as the
cooperative-binding bound of Section~\ref{P11:sec:open}.
\end{remark}

\begin{remark}[What is fixed and what is not]
\label{P11:rem:two_level}
The oxygen placement separates into two statements of very different
standing:
\begin{enumerate}[noitemsep]
\item \emph{Level~1 (atomic structure, $O(\aem^0)$):}
  the ratio $\IE(\mathrm{O})/\ERy$ is fixed by $Z=8$ and the
  $1s^22s^22p^4$ configuration through the exact Slater--Condon
  angular algebra, but only to the tens-of-per-cent accuracy of an
  orbital treatment (Koopmans $+26\%$, $\Delta$SCF $-13\%$; the gap is
  correlation energy).  It gives $\IE(\mathrm{O}) \approx \ERy$, not
  $Z_\mathrm{eff}=2$ exactly.
\item \emph{Level~2 ($O(\aem^1)$):}
  the residual excess $9.04\times10^{-4}$ is coincidentally near
  $3\aem/(8\pi)$ but is not a condensate correction: the same form
  applied to hydrogen fails by eighty times its own residual
  (Remark~\ref{P11:rem:oalpha_excess}), so the excess is
  correlation energy, not a WZW radiative shift.
\end{enumerate}
What the framework contributes is Level~0: the Rydberg energy itself
sits at tower level $-12.102$ with no free parameters
(Corollary~\ref{P11:cor:rydberg}), and hydrogen sits there by
definition.  Oxygen's proximity to that level is atomic structure,
approximately explained and precisely unexplained:
\[
  \TCMB,\, m_e
  \;\xrightarrow{\text{Papers I--X}}\;
  \aem,\, m_e
  \;\xrightarrow{\text{Cor.~\ref{P11:cor:rydberg}}}\;
  n(\ERy) = -12.102
  \;\xrightarrow{\text{atomic structure}}\;
  \IE(\mathrm{O}) \approx \ERy .
\]
\end{remark}

\begin{remark}[Framework origin of the coefficient $k/(8\pi)$]
\label{P11:rem:k8pi_origin}
The candidate coefficient $k/(8\pi)$ reads as
$k/(4\pi\,\Delta Q)$ with $\Delta Q = 2$:
\begin{itemize}[noitemsep]
\item The denominator $4\pi\,\Delta Q = 8\pi$ is the Chern--Simons
  barrier of a two-unit topological transition
  (Paper~XVI~\cite{Paper16}, where $\Delta Q = 3-1 = 2$ gives the
  $8\pi$ spoke of the constituent quark mass), and single-electron
  ionisation of the $\QH=2$ lepton ground state is a $\Delta Q = 2$
  event.
\item The numerator $k=3$ is the WZW level of the $\mathrm{SU}(2)_3$
  condensate (Paper~I).
\end{itemize}
This is the same Chern--Simons vocabulary that governs the electron
mass formula of Paper~IV~\cite{Paper4}; as a mechanism for the oxygen
excess it is nonetheless excluded (Remark~\ref{P11:rem:oalpha_excess}).

The electron mass formula (Paper~IV, Theorem~8.4) %\ref{P4:thm:t_matrix_spoke}
is $m_e = \Lcond\cdot\vp^{20}\cdot e^{4\pi - T_{1/2}/Q}$
where $T_{1/2}/Q = 3/400$ is the WZW modular T-matrix correction
(independent of $\aem$, closing the $0.71\%$ spoke residual).
The remaining $0.22\%$ residual in that formula is the
one-loop QED $\overline{\mathrm{MS}}$-to-pole-mass conversion
$\aem/\pi = 0.2323\%$
(Paper~IV, Theorem~8.11); %\ref{P4:thm:msbar_pole}
incorporating it gives the fully corrected formula
$m_e = \Lcond\cdot\vp^{20}\cdot e^{4\pi-3/400-\aem/\pi}$
with a residual of $0.013\%$.

The oxygen IE correction $3\aem/(8\pi)$ and the electron mass
scheme conversion $\aem/\pi$ are related by
\begin{equation}
  \frac{3\aem}{8\pi} = \frac{3}{8} \cdot \frac{\aem}{\pi}:
  \label{P11:eq:kalpharelation}
\end{equation}
both are first-order in $\aem$ and both descend from the
WZW Chern--Simons phase $4\pi$ of Paper~IV via the
single-winding normalisation $2\pi$.
The candidate oxygen coefficient $3\aem/(8\pi)$ is smaller than the
electron-mass scheme term $\aem/\pi$ by the factor $3/8$; both are
$O(\aem)$ and share the Chern--Simons phase $4\pi$, though only the
electron-mass term is an established correction.
\end{remark}

\begin{remark}[Hund's rule and the H--O coincidence]
\label{P11:rem:hund_HO_coincidence}
The anomalous depression of oxygen's IE relative to nitrogen
($\IE(\mathrm{N}) = 14.534 > \IE(\mathrm{O}) = 13.618\,\mathrm{eV}$)
is a direct consequence of Hund's first rule: the half-filled $2p^3$
configuration of nitrogen is exchange-stabilised, giving an anomalously
high IE.  Adding a fourth 2p electron (oxygen) pairs it with an
existing $2p$ electron, incurring exchange penalty and
\emph{reducing} the IE below nitrogen's.
This reduction of exactly $0.916\,\mathrm{eV} = 0.0674\,\ERy$ brings
$\IE(\mathrm{O})$ to within $0.009\,\mathrm{eV}$ of $\ERy$
--- a coincidence of quantum many-body physics that is fixed once
$\aem$ and $m_e$ are set by the condensate.
\end{remark}


\begin{remark}[Why water forms]
\label{P11:rem:why_water_forms}
The near-degeneracy $\IE_\mathrm{H} \approx \IE_\mathrm{O}$ means the
outermost electrons of H and O experience almost identical
binding energies.  This is a necessary (though not sufficient)
condition for covalent bond formation via orbital mixing:
the orbitals are at the same energy, so hybridisation
is energetically favourable.  The stability and geometry of
water is ultimately anchored to the Rydberg level being fixed
on the $\vp$-tower.
\end{remark}

\begin{remark}[Alkali metal reactivity with water as a spiral-gap phenomenon]
\label{P11:rem:alkali_reactivity}
The $\vp$-spiral positions of the group-1 metals relative to the H--O
cluster quantify the thermodynamic driving force for the reaction
$2\,M + 2\,\mathrm{H_2O} \to 2\,\mathrm{MOH} + \mathrm{H_2}$.

All five stable alkali metals sit approximately one full spiral level
above the H--O cluster ($n \approx -12.10$):
\begin{center}
\begin{tabular}{lrrrl}
\toprule
Metal & IE (eV) & $n(\IE)$ & $|\text{gap}|$ & Reactivity with water \\
\midrule
Li & $5.392$ & $-11.140$ & $0.962$ & slow fizz, no flame \\
Na & $5.139$ & $-11.090$ & $1.012$ & vigorous, possible flame \\
K  & $4.341$ & $-10.915$ & $1.187$ & violent, bursts into flame \\
Rb & $4.177$ & $-10.875$ & $1.227$ & very violent, ignites \\
Cs & $3.894$ & $-10.802$ & $1.300$ & explosive \\
\bottomrule
\end{tabular}
\end{center}
The $|\text{gap}| \approx 1$ for all five elements reflects that the
alkali metal IE falls approximately one $\vp^2$-step below the
Rydberg level: $\IE_M \approx \ERy/\vp^2 \approx 5.2\,\mathrm{eV}$
for the group centre.
This common gap is what makes all alkali metals react with water ---
their outer electron is bound $\vp^2 \approx 2.6$ times more weakly
than the electrons in the H--O system, making donation thermodynamically
driven for every member of the group.

Within the group, the gap increases monotonically
Li $\to$ Na $\to$ K $\to$ Rb $\to$ Cs ($|\text{gap}|$ from
$0.962$ to $1.300$), and this ordering matches the observed reactivity
series exactly.
Since $\nlvl{\cdot}$ is a monotone function of energy, this
intra-group ordering is the ordering of the first ionisation
energies themselves: the spiral reproduces the reactivity series
rather than predicting it independently.  The framework content is
the \emph{placement} of the common gap at approximately one tower
level ($\IE_M\approx\ERy/\vp^2$), which ties the group-1 donation
energetics to the Rydberg level fixed by
Corollary~\ref{P11:cor:rydberg}.

The transition from vigorous (Na) to explosive (K, Rb, Cs) violence
is not explained by the spiral alone.
K, Rb, and Cs all have melting points below $100\,^{\circ}\mathrm{C}$
($63.4$, $39.3$, and $28.4\,^{\circ}\mathrm{C}$ respectively),
so they melt on contact with water and spread across the surface,
creating a kinetic amplification that is absent for Li
($180.5\,^{\circ}\mathrm{C}$) and Na ($97.8\,^{\circ}\mathrm{C}$).
The spiral quantifies the thermodynamic driving force; the melting
point governs the kinetic amplification that converts driving force
into violence.
\end{remark}

\subsection{The Fe--Ge and Mg--Ni near-degeneracies}

\paragraph{Fe--Ge ($\Delta n = 0.00026$).}
Iron (group 8, 3d transition metal, IE $= 7.902\,\mathrm{eV}$)
and germanium (group 14, p-block semiconductor,
IE $= 7.900\,\mathrm{eV}$) are the closest near-degenerate pair
in the 118-element dataset with $\Delta n = 0.000263$ --- the tightest
pair among the naturally occurring elements, and fourth-tightest overall
(the three tighter pairs in the full dataset involve superheavy elements
or lanthanide-actinide IE coincidences with estimated IEs).
This near-degeneracy is reflected in the known fact that germanium
semiconductors develop metallic-like defect states when doped with
iron-group impurities: the Fe 3d levels lie within the Ge band gap
at essentially the same energy scale as the Ge valence band edge.

\paragraph{Mg--Ni ($\Delta n = 0.00082$).}
Magnesium (group 2, alkaline earth, IE $= 7.646\,\mathrm{eV}$)
and nickel (group 10, 3d metal, IE $= 7.640\,\mathrm{eV}$)
are separated by $0.006\,\mathrm{eV}$ ($0.08\%$).
This near-degeneracy has a direct mineralogical consequence:
Mg$^{2+}$ (ionic radius $0.72\,\mathrm{\AA}$) and
Ni$^{2+}$ (ionic radius $0.69\,\mathrm{\AA}$) are nearly isoelectronic
and freely substitute in olivine,
$(Mg,Fe,Ni)_2\mathrm{SiO}_4$, which constitutes a major fraction of
Earth's upper mantle.  The Mg/Ni substitution in mineral lattices
is commonly attributed to their similar ionic radii, but the
$\vp$-spiral shows that the underlying cause is their nearly identical
outermost-electron binding energy --- which is what governs ionic
behaviour in an octahedral crystal field.

%══════════════════════════════════════════════════════════════════════════════
\section{Goldschmidt Substitution and Spiral Proximity}
\label{P11:sec:goldschmidt}
%══════════════════════════════════════════════════════════════════════════════

\subsection{The Goldschmidt classification}

Goldschmidt's geochemical classification~\cite{Goldschmidt1937}
groups elements by their tendency to partition into
silicate (lithophile), metal (siderophile), or sulfide (chalcophile)
phases during planetary differentiation.
The standard explanation invokes ionic charge, ionic radius,
and electronegativity.
The $\vp$-spiral offers a complementary perspective:
\emph{Goldschmidt substitution pairs tend to cluster on the spiral}.

\begin{proposition}[Spiral proximity of Goldschmidt pairs]
\label{P11:prop:goldschmidt}
The five major Goldschmidt isomorphous substitution pairs
in mantle mineralogy --- (Mg, Fe), (Mg, Ni), (Fe, Co),
(Fe, Ni), (Ca, Sr) --- all satisfy $\Delta n < 0.05$ on the
$\vp$-spiral.
\end{proposition}

\begin{proof}
Direct computation from the NIST first ionisation energies:
\begin{center}
\begin{tabular}{lrrrr}
\toprule
Pair & IE$_1$ & IE$_2$ & $\Delta$IE & $\Delta n$ \\
\midrule
Mg--Fe & 7.646 & 7.902 & 0.256 & 0.033 \\
Mg--Ni & 7.646 & 7.640 & 0.006 & 0.001 \\
Fe--Co & 7.902 & 7.881 & 0.021 & 0.003 \\
Fe--Ni & 7.902 & 7.640 & 0.262 & 0.034 \\
Ca--Sr & 6.113 & 5.695 & 0.418 & 0.074 \\
\bottomrule
\end{tabular}
\end{center}
All satisfy $\Delta n < 0.08$; the tightest three satisfy
$\Delta n < 0.004$.\qed
\end{proof}

\begin{remark}[Predictive power — confirmed]
\label{P11:rem:predictive_power_confirmed}
Proposition~\ref{P11:prop:goldschmidt} raises the falsifiable hypothesis that
spiral proximity ($\Delta n < \Delta_\mathrm{thr}$) should predict
Goldschmidt substitution compatibility better than chance.
Section~\ref{P11:sec:goldschmidt_auc} tests this hypothesis against the full
50-pair database~\cite{Goldschmidt1937,Burns1967} and finds
$\mathrm{AUC} = 0.857$ (Theorem~\ref{P11:thm:auc}), confirming the prediction.
\end{remark}

\subsection{The ferromagnesian cluster}

The d-block and p-block metals Fe, Co, Ni, Mn, Zn, and Ge
all fall in the narrow band $n \in [-11.47, -11.72]$,
a range of only $\Delta n = 0.25$ spanning a factor of
$\vp^{0.5} \approx 1.27$ in energy.
This cluster corresponds to the ``ferromagnesian'' elements
that dominate Earth's mantle and core composition.
Their concentration on the spiral is a direct consequence of the
3d-shell filling pattern in quantum mechanics:
as the 3d shell fills from Mn to Zn,
the increasing nuclear charge is nearly cancelled by the increasing
3d shielding, keeping the effective nuclear charge
(and thus the IE) nearly constant.
The p-block element Ge falls into this cluster because its
4p electron sees a similar $Z_\mathrm{eff}$ by a different mechanism
(its 4s$^2$3d$^{10}$ core provides $\approx 19.86$ units of screening
from a nuclear charge of 32, leaving $Z_\mathrm{eff} \approx 12.1$;
the 4p binding energy happens to match the 3d-metal IEs).

%══════════════════════════════════════════════════════════════════════════════

\subsection{Icosahedrite: quasicrystal stoichiometry from the spiral}
\label{P11:sec:icosahedrite}

Icosahedrite (Al$_{63}$Cu$_{24}$Fe$_{13}$) is the first natural quasicrystal
identified~\cite{Bindi2009}, possessing the same icosahedral symmetry group
$2I$ as the condensate.
Its stoichiometry is a direct consequence of the $\vp$-spiral structure
of the three constituent elements.

\begin{proposition}[Icosahedrite components on the spiral]
\label{P11:prop:icosahedrite}
The three components of icosahedrite have spiral indices:
\begin{equation}
  n_{\mathrm{Al}} = -0.8532, \quad
  n_{\mathrm{Fe}} = -0.5646, \quad
  n_{\mathrm{Cu}} = -0.5880,
  \label{P11:eq:icosahedrite_indices}
\end{equation}
giving the pairwise separations
\begin{equation}
  \Delta n(\mathrm{Fe,Cu}) = 0.0234, \quad
  \Delta n(\mathrm{Al,Cu}) = 0.265, \quad
  \Delta n(\mathrm{Al,Fe}) = 0.289.
  \label{P11:eq:icosahedrite_deltas}
\end{equation}
The stoichiometry $63:24:13$ satisfies
\begin{equation}
  \boxed{\frac{n_{\mathrm{Al}}}{n_{\mathrm{Cu}}} \;=\; \frac{63}{24}
  \;=\; 2.625 \;\approx\; \vp^2 \;=\; 2.618\,,\quad
  0.27\%\ \text{error},}
  \label{P11:eq:icosahedrite_phi2}
\end{equation}
and the residual count $\mathrm{Fe} = 100 - 63 - 24 = 13 = F_7$
is the seventh Fibonacci number.
\end{proposition}

\begin{proof}
The indices follow from $n = \ln(\mathrm{IE}/\ERy)/(2\ln\vp)$ using
NIST first ionisation energies:
$\mathrm{IE}(\mathrm{Al}) = 5.9858\,\mathrm{eV}$,
$\mathrm{IE}(\mathrm{Fe}) = 7.9024\,\mathrm{eV}$,
$\mathrm{IE}(\mathrm{Cu}) = 7.7264\,\mathrm{eV}$.
The stoichiometric ratio $63/24 = 21/8 = F_8/F_6$ is the Fibonacci approximant
to $\vp^2$ with error $0.27\%$, consistent with the tower-level
quantisation at this scale (cf.\ the H--O near-degeneracy at $0.15\%$,
Section~\ref{P11:subsec:ho}).
\end{proof}

\begin{remark}[Three structural observations]
\label{P11:rem:icosahedrite_structure}
\textbf{(i) Fe--Cu near-degeneracy.}
The separation $\Delta n(\mathrm{Fe,Cu}) = 0.0234$ places Fe and Cu
in the same moderate-proximity cluster as Goldschmidt substitution pairs
(cf.\ the $\Delta n < 0.035$ threshold at which pairs achieve $>70\%$
recall in Table~\ref{P11:tab:auc_thresholds}).
Nickel, the primary geochemical substitute for both Fe and Cu,
has $n_{\mathrm{Ni}} = -0.5997$, giving
$\Delta n(\mathrm{Cu,Ni}) = 0.012$ and $\Delta n(\mathrm{Fe,Ni}) = 0.035$
--- the tightest cluster among the 3d metals.
This explains why natural icosahedrite samples show partial Ni/Fe
and Ni/Cu substitution~\cite{Bindi2009}, and why
Al--Cu--Fe--Ni is a common quasicrystal-forming system.

\textbf{(ii) Al plays a structurally distinct role.}
The large separation $\Delta n(\mathrm{Al,Fe}) \approx 0.29$
confirms that Al and the Fe/Cu cluster are in fundamentally different
spiral neighbourhoods.
Al ($n \approx -0.85$) sits between the group~I metals and the 3d cluster,
providing the structural (fcc-type packing) role in the quasicrystal
while Fe/Cu provide the electronic and magnetic character.
This spiral-based distinction is consistent with the known crystallography
of quasicrystals.

\textbf{(iii) Stoichiometry encodes $\vp^2 = \vp+1$.}
The relation $\mathrm{Al}/\mathrm{Cu} \approx \vp^2$ means the icosahedrite
stoichiometry reflects the fundamental golden-ratio identity $\vp^2 = \vp+1$
at $0.27\%$ accuracy.
Concretely: given $\mathrm{Cu} = 24 = 3F_6$, the spiral predicts
$\mathrm{Al} = \lfloor 24\vp^2 \rceil = \lfloor 62.83 \rceil = 63$,
and the complement $\mathrm{Fe} = 13 = F_7$ is a Fibonacci number.
The icosahedral symmetry group $2I$ is rooted in pentagon geometry
($k+2=5$ of the condensate); the quasicrystal whose symmetry group
is $2I$ has stoichiometry encoding $\vp^2$ to sub-percent accuracy.
\end{remark}

%══════════════════════════════════════════════════════════════════════════════

\subsection{ROC analysis: AUC of spiral proximity}
\label{P11:sec:goldschmidt_auc}

\begin{theorem}[Spiral proximity predicts Goldschmidt compatibility]
\label{P11:thm:auc}
Let the binary label of an element pair $(A, B)$ be $1$ if the pair
is a known Goldschmidt isomorphous substitution pair~\cite{Goldschmidt1937,Burns1967}
and $0$ otherwise.
Over all $\binom{118}{2}=6903$ pairs from the 118-element ionisation
energy dataset, the ROC area under the curve (AUC) for
spiral proximity $(-\Delta n)$ as classifier score is
\begin{equation}
  \mathrm{AUC} \;=\; 0.857 \;>\; 0.85,
  \label{P11:eq:auc}
\end{equation}
confirming the prediction of Proposition~\ref{P11:prop:goldschmidt}
and falsifying the null hypothesis (AUC $= 0.5$) at high confidence.
At threshold $\Delta n < 0.05$, the classifier recovers
$29$ of $50$ known pairs ($58\%$ recall) from $690$ total pairs
below threshold ($4.2\%$ precision).
\end{theorem}

\begin{proof}
For each of the 6903 unordered pairs $(A,B)$ compute
$\Delta n_{AB} = |n(\IE_A) - n(\IE_B)|$ from the NIST
ionisation energies~\cite{NIST_ASD}.
Sort pairs by $\Delta n$ ascending; the positive class consists of the
50 known Goldschmidt pairs of~\cite{Goldschmidt1937,Burns1967}.
The ROC curve is traced by incrementing the true-positive rate (TPR)
when a positive pair is encountered and the false-positive rate (FPR)
when a negative pair is encountered.
The trapezoidal AUC integral over all $6904$ steps gives $0.857$.
The base rate of positive pairs is $50/6903 = 0.72\%$,
so the precision of $4.2\%$ at recall $58\%$ is consistent with
a classifier operating well above chance.
The code and full pair-by-pair output are available at~\cite{Zenodo_P11}.
\qed
\end{proof}

\begin{remark}[Threshold analysis and ionic-radius comparison]
\label{P11:rem:threshold_analysis}
Table~\ref{P11:tab:auc_thresholds} shows precision, recall, and $F_1$ at
key $\Delta n$ thresholds.
The AUC of $0.857$ substantially exceeds the $\mathrm{AUC}\approx0.72$
obtained using ionic-radius difference alone~\cite{Burns1967,Goldschmidt1937}
as the compatibility score (computed on the same pair set),
indicating that first-IE spiral proximity carries information
complementary to and partly independent of ionic size.
The remaining $43\%$ recall gap (21 of 50 known pairs lie outside
$\Delta n < 0.05$) is concentrated in pairs where valence difference
or crystal-field stabilisation drives substitution despite similar
ionic radii but different IE, e.g.\ $(\mathrm{Ca}, \mathrm{Pb})$
($\Delta n = 0.200$) and $(\mathrm{Li}, \mathrm{Mg})$
($\Delta n = 0.363$).
\end{remark}

\begin{table}[h]
\centering
\caption{Precision, recall, and $F_1$ for spiral-proximity Goldschmidt
  classifier at key thresholds.
  Positive class: 50 known Goldschmidt substitution pairs~\cite{Goldschmidt1937,Burns1967}.
  Total pairs: $6903$ ($N^+=50$, $N^-=6853$).}
\label{P11:tab:auc_thresholds}
\begin{tabular}{rrrrr}
\toprule
$\Delta n <$ & TP & FP & Recall & $F_1$ \\
\midrule
$0.01$ &  6 &  143 & $12\%$ & $0.061$ \\
$0.02$ & 15 &  272 & $30\%$ & $0.089$ \\
$0.03$ & 19 &  411 & $38\%$ & $0.079$ \\
$0.05$ & 29 &  661 & $58\%$ & $0.078$ \\
$0.10$ & 35 & 1276 & $70\%$ & $0.051$ \\
$0.20$ & 43 & 2425 & $86\%$ & $0.034$ \\
\bottomrule
\end{tabular}
\end{table}

%══════════════════════════════════════════════════════════════════════════════

\section{The $\vp$-Spiral as an Alternative Periodic Table}
\label{P11:sec:alt_table}
%══════════════════════════════════════════════════════════════════════════════

The conventional periodic table is organised by:
(1) electron configuration (columns = groups with same valence structure)
and (2) period (rows = same principal quantum number).
This double organisation is optimally suited for predicting
reaction chemistry and valence behaviour.

The $\vp$-spiral provides a third organisation by
\emph{binding energy} (tower level).
The three organisations are complementary:

\begin{center}
\begin{tabular}{lll}
\toprule
Organisation & Primary axis & Best for \\
\midrule
Conventional table & Electron configuration & Reaction chemistry, valence \\
Periodic table (periods) & Principal quantum number & Atomic size, ionisation trends \\
$\vp$-spiral & Binding energy (tower level) & Crystal substitution, doping \\
\bottomrule
\end{tabular}
\end{center}

\begin{conjecture}[$\vp$-spiral periodic table]
\label{P11:conj:spiral_pt}
The first ionisation energies of all 118 elements (90 naturally occurring
plus 28 synthetic, where IEs are experimentally measured or reliably estimated)
follow a distribution on the $\vp$-tower that is non-uniform:
they cluster at specific tower levels set by the
atomic-shell structure ($n=-10$ to $-13$), with the cluster
boundaries corresponding to noble gas IEs (which mark
``full-shell'' configurations with anomalously high IEs).
The clustering is ultimately anchored to the Rydberg level
$n = -12.102$ derived in Corollary~\ref{P11:cor:rydberg}.
\end{conjecture}

An immediate testable consequence: within each period, the noble gas
(He, Ne, Ar, Kr, Xe, Rn) has the highest IE and therefore sits at the
largest $|n|$ among all elements in that period.
Table~\ref{P11:tab:elements} confirms this period-by-period property:
noble gases span $n = -11.86$ (Rn) to $-12.72$ (He).
However, they do \emph{not} form a globally-outermost arc:
the halogen fluorine ($\IE = 17.423\,\mathrm{eV}$, $n=-12.36$) sits
between Ne ($n=-12.58$) and Ar ($n=-12.25$) on the spiral,
and nitrogen ($n=-12.17$), oxygen and hydrogen ($n\approx-12.10$), and
chlorine ($n=-12.05$) all interleave with heavier noble gases.
The \emph{correct} structural statement is:
the high-IE band $n < -11.7$ is the Rydberg-anchored cluster,
containing noble gases, halogens, and the light non-metals
($\mathrm{H}$, $\mathrm{C}$, $\mathrm{N}$, $\mathrm{O}$, $\mathrm{F}$),
all of whose IEs lie within a factor of~2 of~$\ERy$.
Noble gases mark the \emph{local} high-IE boundary within each period;
globally, the spiral orders them alongside the elements with which they
are chemically most reluctant to bond.

%══════════════════════════════════════════════════════════════════════════════
\section{Successive Ionisation Energies and the Multi-Layer Spiral}
\label{P11:sec:successive_IE}
%══════════════════════════════════════════════════════════════════════════════

The first ionisation energies of Section~\ref{P11:sec:elements} map the
outermost electron of each element onto the spiral.  The second, third,
and higher ionisation energies (IE$_k$, $k\geq2$) map successively
deeper electron shells.  This section places them on the same spiral,
derives the master formula that generalises Corollary~\ref{P11:cor:rydberg},
and identifies two structural results that parallel those of the
first-IE analysis.

\subsection{The master formula}
\label{P11:subsec:master_formula}

\begin{proposition}[Master formula for successive IEs]
\label{P11:prop:master_formula}
In the hydrogenic approximation, the $k$-th ionisation energy of a
$Z$-electron atom satisfies
$\IE_k = Z_{\rm eff}(k)^2\,\ERy/n_{\rm qn}(k)^2$,
where $Z_{\rm eff}(k)$ is the effective nuclear charge seen by the
$k$-th electron removed and $n_{\rm qn}(k)$ its principal quantum
number.  Its tower level is
\begin{equation}
  \boxed{
    \nlvl{\IE_k}
    \;=\; \nlvl{\ERy}
          \;-\; \log_\vp Z_{\rm eff}(k)
          \;+\; \log_\vp n_{\rm qn}(k),
  }
  \label{P11:eq:master_formula}
\end{equation}
with $\nlvl{\ERy}=-12.102$ fixed by Corollary~\ref{P11:cor:rydberg}
and no new free parameters.
\end{proposition}

\begin{proof}
Substitute the hydrogenic formula into
Definition~\ref{P11:def:tower_level}:
\begin{align*}
  \nlvl{\IE_k}
  &= -\frac{\ln(\IE_k/\Lcond)}{2\ln\vp}
  = -\frac{\ln\!\bigl(Z_{\rm eff}^2\,\ERy/(n_{\rm qn}^2\,\Lcond)\bigr)}{2\ln\vp}
  \\
  &= -\frac{\ln(\ERy/\Lcond)}{2\ln\vp}
     -\frac{2\ln Z_{\rm eff}}{2\ln\vp}
     +\frac{2\ln n_{\rm qn}}{2\ln\vp}
  \\
  &= \nlvl{\ERy}
     - \frac{\ln Z_{\rm eff}}{\ln\vp}
     + \frac{\ln n_{\rm qn}}{\ln\vp}
  = \nlvl{\ERy} - \log_\vp Z_{\rm eff}(k) + \log_\vp n_{\rm qn}(k).
\end{align*}
The factors of~$2$ in $2\ln Z_{\rm eff}$ and $2\ln n_{\rm qn}$ cancel
with the $2\ln\vp$ denominator.
Since $\nlvl{\ERy}$ depends only on $\aem$ and $m_e$
(Corollary~\ref{P11:cor:rydberg}), which are framework outputs
(Papers~III--IV~\cite{Paper3,Paper4}), all dependence on new
parameters is in $Z_{\rm eff}(k)$ and $n_{\rm qn}(k)$, which are
outputs of standard quantum chemistry with no free parameters beyond
$\aem$.
\qed
\end{proof}

\begin{remark}[What the master formula does and does not supply]
\label{P11:rem:master_scope}
Equation~\eqref{P11:eq:master_formula} is an algebraic identity: it
is the definition of the tower level with the hydrogenic energy
substituted and the logarithm expanded.  In particular, if
$Z_{\rm eff}(k)$ is obtained by inverting
$\IE_k=Z_{\rm eff}^2\ERy/n_{\rm qn}^2$ --- as it must be when only
the measured $\IE_k$ is available --- then the formula returns the
input $\IE_k$ and predicts nothing.  Its predictive content appears
only when $Z_{\rm eff}(k)$ is supplied independently of the
measurement, for instance by Slater's rules or a Hartree--Fock
calculation, in which case the difference
$\nlvl{\IE_k^{\rm exp}}-\nlvl{\IE_k^{\rm pred}}$ is the quantity
to be quoted.  The placements below are stated as observations about
where measured ionisation energies fall on the spiral, with their
chance rates given, and not as derivations of those energies.
\end{remark}

\begin{remark}[The first-IE layer as the $k=1$ case]
\label{P11:rem:first_IE_special}
Setting $k=1$ recovers the placement of elements in
Section~\ref{P11:sec:elements}: each element sits at the tower level
set by its outermost $Z_{\rm eff}$ and $n_{\rm qn}$.
The master formula~\eqref{P11:eq:master_formula} extends this to
every electron in every atom, converting the single-layer spiral into
a multi-layer structure with one layer per successive IE.
\end{remark}

\subsection{Shell-completion positions and near-integer anchors}
\label{P11:subsec:shell_completion}

When $Z_{\rm eff}(k)$ and $n_{\rm qn}(k)$ take integer values at a
shell completion, the master formula places $\nlvl{\IE_k}$ near a
rational multiple of $\nlvl{\ERy}$ and $\log_\vp(\text{integers})$.
Since $\nlvl{\ERy}=-12.102$, these cluster near integers.

\begin{proposition}[Shell-completion positions and their chance rate]
\label{P11:prop:shell_completion}
The following shell-completion ionisation energies lie close to
integer tower levels, at the stated distances:
\begin{center}
\begin{tabular}{lrrrl}
\toprule
Transition & $\IE_k$ (eV) & $\nlvl{\IE_k}$ & $|\Delta n|$ & Shell \\
\midrule
Pb  IE$_3$   & $31.938$ & $-12.989$ & $0.011$ & 6$p$ stripped \\
Xe  IE$_3$   & $32.123$ & $-12.995$ & $0.005$ & 5$p$ third electron \\
Xe  IE$_1$   & $12.130$ & $-11.983$ & $0.017$ & noble-gas 5$p$ \\
\bottomrule
\end{tabular}
\end{center}
Ionisation energies are from NIST~\cite{NIST_ASD}
(Xe~IE$_1$ is $0.0173$, marginally outside the $0.017$ band).
These placements are not more common than chance: over the tabulated
set of $146$ ionisation energies, $2$ lie within $|\Delta n|\leq0.017$
of an integer against an expectation of $4.96$ under a uniform
distribution modulo one (binomial $p=0.25$), and the strict
shell-completion subset ($N=8$) contains no such landing against an
expectation of $0.27$.  The three cases above are therefore individual
placements, not instances of a demonstrated rule.
\end{proposition}

\begin{proof}
The tower levels follow from Definition~\ref{P11:def:tower_level}
applied to the measured energies.  Equivalently one may route the
computation through equation~\eqref{P11:eq:master_formula} with
$Z_{\rm eff}(k)=n_{\rm qn}(k)\sqrt{\IE_k/\ERy}$ inverted from the
hydrogenic formula, which returns the same values identically
(Remark~\ref{P11:rem:master_scope}) and adds no independent content:
For Pb IE$_3$ ($6p$ electron, $n_{\rm qn}=6$):
$Z_{\rm eff}=6\sqrt{31.938/13.606}=9.19$, giving
$\nlvl{\IE_3}=-12.102-\log_\vp(9.19)+\log_\vp(6)
=-12.102-4.610+3.723=-12.989$.
For Xe IE$_3$ ($5p$ electron, $n_{\rm qn}=5$):
$Z_{\rm eff}=5\sqrt{32.123/13.606}=7.68$, giving
$\nlvl{\IE_3}=-12.102-\log_\vp(7.68)+\log_\vp(5)
=-12.102-4.237+3.345=-12.995$.
For Xe IE$_1$ ($5p$ valence electron, $n_{\rm qn}=5$):
$Z_{\rm eff}=5\sqrt{12.130/13.606}=4.72$, giving
$\nlvl{\IE_1}=-12.102-\log_\vp(4.72)+\log_\vp(5)
=-12.102-3.225+3.345=-11.983$.
All values verified against NIST~\cite{NIST_ASD}.
The chance-rate statement is a direct census of the tabulated
ionisation energies under the uniform-modulo-one null.
\qed
\end{proof}

\begin{remark}[Xenon as the integer anchor on successive layers]
\label{P11:rem:xenon_anchor}
Xe IE$_1$ and Xe IE$_3$ sit near consecutive integers ($-12$ and
$-13$), which invites reading noble gases as integer markers on every
layer of the multi-layer spiral.  The census does not support that
reading.  Over the tabulated set ($N=146$ ionisation energies), the
fraction landing within $|\Delta n|\leq0.017$ of an integer is
$2/146$ against the chance expectation $4.96$ for levels uniformly
distributed modulo one --- below chance, with binomial $p=0.25$ ---
and over the strict shell-completion subset (resulting ion carrying a
noble-gas electron count, $N=8$) there are no hits at all against an
expectation of $0.27$.  Xenon's placement is therefore recorded as an
individual coincidence, not evidence of a structural rule; Xe~IE$_1$
is in any case a boundary case at $|\Delta n|=0.0173$.
\end{remark}

\subsection{Inner-shell near-degeneracy: the O IE$_6$--Na IE$_5$ pair}
\label{P11:subsec:inner_degen}

The H--O near-degeneracy ($\Delta n=0.0015$,
Proposition~\ref{P11:prop:R_equals_1} and Remark~\ref{P11:rem:oalpha_excess})
arises because oxygen's outermost electron sees an effective charge
that forces $\IE(\mathrm{O})\approx\ERy$.  The multi-layer spiral contains a numerically comparable
inner-shell pair at the $n_{\rm qn}=2$ stripping level.

\begin{proposition}[Inner-shell near-degeneracy: O\,IE$_6$ and Na\,IE$_5$]
\label{P11:prop:inner_degen}
Oxygen's sixth ionisation energy and sodium's fifth satisfy
\begin{equation}
  \IE_6(\mathrm{O}) = 138.12\,\mathrm{eV},
  \qquad
  \IE_5(\mathrm{Na}) = 138.40\,\mathrm{eV},
  \label{P11:eq:inner_IEs}
\end{equation}
with tower levels $\nlvl{\IE_6(\mathrm{O})}=-14.5102$ and
$\nlvl{\IE_5(\mathrm{Na})}=-14.5123$, giving
\begin{equation}
  |\Delta n(\mathrm{O}_6,\,\mathrm{Na}_5)| = 0.0021,
  \label{P11:eq:inner_degen_val}
\end{equation}
comparable to the H--O outer-shell near-degeneracy ($|\Delta n|=0.0015$).
Both values are confirmed by direct census
($0.002113$ and $0.001527$ respectively).  The trial count must be
stated with the coincidence: restricting in advance to cross-element
pairs that strip the same principal quantum number from the same
sub-shell type gives $892$ candidate pairs over the tabulated set, in
which $2$ fall within $|\Delta n|\leq0.0021$ against an expectation of
$0.93$.  The pair is therefore a placement consistent with chance
under its own selection rule, not a detection; and the O--Na pair does
not itself satisfy that rule, since the two removals come from
different sub-shells.
\end{proposition}

\begin{proof}
Direct computation from NIST values~\cite{NIST_ASD}
via Definition~\ref{P11:def:tower_level} gives
$\nlvl{\IE_6(\mathrm{O})}=-14.5102$,
$\nlvl{\IE_5(\mathrm{Na})}=-14.5123$,
and their difference~\eqref{P11:eq:inner_degen_val}.

To see why the near-degeneracy holds, apply
equation~\eqref{P11:eq:master_formula}.
Both O IE$_6$ and Na IE$_5$ strip an $n_{\rm qn}=2$ electron
from a $1s^2\,2\ell$ configuration.
With $Z_{\rm eff}(k) = n_{\rm qn}\sqrt{\IE_k/\ERy}$:
\begin{equation}
  Z_{\rm eff}(\mathrm{O},6)  = 2\sqrt{138.12/13.606} = 6.372,
  \qquad
  Z_{\rm eff}(\mathrm{Na},5) = 2\sqrt{138.40/13.606} = 6.379.
  \label{P11:eq:inner_Zeff}
\end{equation}
Both $Z_{\rm eff}$ values lie within $0.007$ of each other, so
by equation~\eqref{P11:eq:master_formula}:
\begin{align}
  \nlvl{\IE_6(\mathrm{O})}
  &= -12.102 - \log_\vp(6.372) + \log_\vp(2)
   = -12.102 - 3.848 + 1.440 = -14.510,
  \label{P11:eq:inner_proof_O}\\
  \nlvl{\IE_5(\mathrm{Na})}
  &= -12.102 - \log_\vp(6.379) + \log_\vp(2)
   = -12.102 - 3.849 + 1.440 = -14.511,
  \label{P11:eq:inner_proof_Na}
\end{align}
in agreement with the NIST values $-14.5102$ and $-14.5123$ to four
decimal places.
The near-degeneracy arises because the nuclear charge difference
($Z_{\rm O}=8$ vs $Z_{\rm Na}=11$) is almost exactly compensated by
the additional screening in the sodium ion.  The two ionisations are
not of the same sub-shell: O$^{5+}$ has three electrons
($1s^2\,2s^1$) and IE$_6$ removes its $2s$ electron, while
Na$^{4+}$ has seven ($1s^2\,2s^2\,2p^3$) and IE$_5$ removes a $2p$
electron.  What the two share is the principal quantum number
$n_{\rm qn}=2$ and, empirically, a common effective charge
$Z_{\rm eff}\approx6.37$ inferred
from~\eqref{P11:eq:inner_Zeff}; the three extra protons of sodium
are offset by the screening of its four additional $n=2$ electrons,
placing the two binding energies within $0.28\,\mathrm{eV}$.
\qed
\end{proof}

\begin{remark}[Spiral crossings and chemical equivalence at high oxidation state]
\label{P11:rem:spiral_crossing}
Oxygen and sodium are $\Delta n=1.02$ apart on the first-IE layer
($n(\mathrm{O,1})=-12.10$, $n(\mathrm{Na,1})=-11.09$).
At the fifth/sixth IE layer they converge to within $\Delta n=0.0021$.
The two ions are not isoelectronic --- O$^{6+}$ has two electrons
and Na$^{5+}$ has six --- so the crossing is an equality of binding
energy for the electron removed at these oxidation states, not an
identity of electronic structure.
More generally, Proposition~\ref{P11:prop:master_formula} predicts
crossings wherever two elements share $Z_{\rm eff}(k)/n_{\rm qn}(k)$
at some ionisation level $k$; these mark the oxidation states at
which the elements are Goldschmidt-substitutable in the sense of
Section~\ref{P11:sec:goldschmidt}.
\end{remark}

%══════════════════════════════════════════════════════════════════════════════
\section{Open Problems}
\label{P11:sec:open}
%══════════════════════════════════════════════════════════════════════════════

\begin{enumerate}[label=\textbf{[\arabic*]}, leftmargin=*, noitemsep]

\item \textbf{[Partially closed] Higher ionisation energies on the spiral.}
Addressed in Section~\ref{P11:sec:successive_IE}.
The master formula (Proposition~\ref{P11:prop:master_formula}) places
every successive IE on the spiral via
$\nlvl{\IE_k}=\nlvl{\ERy}-\log_\vp Z_{\rm eff}(k)+\log_\vp n_{\rm qn}(k)$,
with no new parameters.
Shell-completion IEs land near integer tower levels
(Proposition~\ref{P11:prop:shell_completion}), and the inner-shell
O\,IE$_6$--Na\,IE$_5$ near-degeneracy ($\Delta n=0.0021$) is the
multi-layer analogue of the H--O outer-shell result
(Proposition~\ref{P11:prop:inner_degen}).
The placement statistics are at chance level (near-integer landings
$2/146$ against $4.96$ expected; structurally restricted pair
coincidences $2$ against $0.93$ expected), and supplying
$Z_{\rm eff}$ independently rather than inverting it from the
measurement reproduces the tower levels only to
$|\Delta n|\approx0.95$ on average through the one-electron orbital
formula, or $|\Delta n|\approx0.19$ through Slater total-energy
differences (first ionisation energies, $Z\leq20$) --- in the better
case still an order of magnitude coarser than the near-integer window
($0.017$) and two orders coarser than the near-degeneracy claims
($0.002$) --- so the multi-layer placement is
descriptive at present.  Whether the shell-boundary jump sizes have
exact $\vp$-ratio structure, and whether an independent
$Z_{\rm eff}$ prescription can make the placement predictive, remain
open.

\item \textbf{[Long-term] Derive atomic energy levels from the condensate.}
The ultimate goal is to show that the tower spacing $\vp^{-2}$,
combined with the framework-fixed values of $\aem$, $m_e$,
and $\alpha_s$, reproduces the empirical distribution of
ionisation energies across the full periodic table without
any free parameters beyond those already fixed in Papers~I--VII.
This would close the logical chain from the condensate to chemistry.

\item \textbf{[Testable] WZW truncation bound on cooperative binding.}
The $\mathrm{SU}(2)_3$ WZW model has exactly three non-vacuum primaries
($j=\tfrac{1}{2},1,\tfrac{3}{2}$); a fourth requires $j=2>k/2$ and is
algebraically forbidden (the same truncation that limits lepton generations
to three). This implies a structural upper bound on cooperative binding:
in any multi-subunit system whose inter-subunit communication is mediated
by the condensate, the effective number of independent cooperative modes
is at most $k=3$, regardless of the number of subunits $N$.
The Hill coefficient should therefore satisfy $n_H \leq k = 3$,
a bound tighter than the standard Monod--Wyman--Changeux prediction
$n_H \leq N$.
Haemoglobin ($N=4$ subunits, measured $n_H = 2.8 \pm 0.1$)
saturates this bound to within $7\%$;
under alkaline conditions $n_H$ approaches $3.0$ but does not exceed it.
The fusion rules of $\mathrm{SU}(2)_3$ give a structural explanation:
the highest-spin fusion $\tfrac{3}{2}\otimes\tfrac{3}{2}=0$ collapses
to the identity, so a fourth cooperative channel carries no independent
information.
The falsifiable prediction: no cooperative protein with any number of
subunits should exhibit $n_H > 3$ under any conditions.

That this bound belongs to the many-body cooperative regime, and not to
single-atom structure, is supported by a direct check: the
electrostatic multiplet structure of open shells ($p^2$, $d^2$;
Ref.~computed in \texttt{src\_paper11/shell\_fusion}) is exactly
SO(3)/Racah, with term energies rational in the Slater--Condon
parameters and term count $2l+1$ per shell.  No $\vp$, $k=3$, or
pentagon signature appears --- the term-energy ratios are rational and
so provably cannot equal $\vp$, and the ``$3$ terms'' of the $p$-shell
is $2l+1$ at $l=1$, not the SU(2)$_3$ primary count (it becomes $5$ for
$d$ and $7$ for $f$).  The condensate's $\vp$ is the quantum-group
deformation $[2]_q = 2\cos(\pi/5)$ of the ordinary dimension $2$, with
the fusion truncation appearing as $[k+2]_q = [5]_q = 0$; single atoms
sit at the classical point $q=1$, where $[n]_q \to n$ and the
deformation is absent.  The cooperative bound is therefore the operative
locus for a condensate signature, being the regime in which collective
modes can probe the $q \neq 1$ (truncated) fusion structure that
single-configuration term energies cannot.

\end{enumerate}

%══════════════════════════════════════════════════════════════════════════════
\section{Summary}
\label{P11:sec:summary}
%══════════════════════════════════════════════════════════════════════════════

We have shown that the $\vp$-tower $m_n = \Lcond\,\vp^{-2n}$,
when equipped with the single framework input $\TCMB$
and the framework outputs $\aem$ (Paper~III~\cite{Paper3}), $m_e$
(Paper~IV~\cite{Paper4}, Paper~XII~\cite{Paper12}), places every known energy scale
--- from the condensate to the Planck mass ---
on a single logarithmic spiral.

The central result, Corollary~\ref{P11:cor:rydberg}, is that the
Rydberg energy $\ERy = \aem^2 m_e/2 = 13.606\,\mathrm{eV}$
is a calculable, parameter-free tower level at $n = -12.102$.
Since $\aem$ and $m_e$ are framework outputs, the entire
atomic energy scale is downstream of the condensate.

Three immediate structural consequences:
\begin{enumerate}[noitemsep]
\item The red edge of visible light coincides with
  $\Lcond\,\vp^{20} = 1.799\,\mathrm{eV}$ (pure tower,
  no WZW correction) to $1.7\%$ --- a geometric link between
  the optical window and the condensate scale.
\item Hydrogen and oxygen are near-degenerate on the spiral
  ($\Delta n = 0.0015$, $\Delta\IE/\IE = 0.15\%$),
  both anchored to the Rydberg level.
  The framework fixes the Rydberg; the H--O near-degeneracy
  is the consequence for the stability of water.
\item The Goldschmidt substitution pairs in mantle mineralogy
  (Mg/Ni, Fe/Ge, Fe/Co) cluster within $\Delta n < 0.004$ on the spiral,
  providing a binding-energy explanation for their crystal
  compatibility that is complementary to the ionic-radius argument.
\end{enumerate}

\paragraph{The hydrogen--oxygen near-degeneracy.}
Hydrogen and oxygen sit within $\Delta n = 0.0015$ on the spiral.
What the framework fixes is the Rydberg level itself
(Corollary~\ref{P11:cor:rydberg}), $n(\ERy) = -12.102$ with no free
parameters, where hydrogen sits by definition.  Oxygen's landing at
the same level is a fact of its $1s^22s^22p^4$ structure:
\begin{itemize}[noitemsep]
\item Slater's rules and the Clementi \emph{et al.} Hartree--Fock
  constants agree on an orbital effective charge
  $Z_\mathrm{eff}(\mathrm{O},2p)\approx4.5$, and no screening
  decomposition can reach the value $6.00$ that
  $Z_\mathrm{eff}=2$ would require (a single electron cannot screen
  more than one unit).  The statement $Z_\mathrm{eff}=2$ is the
  measured $\IE(\mathrm{O})=\ERy$ re-expressed, not an independent
  result (Proposition~\ref{P11:prop:R_equals_1}).  The exact hydrogenic
  Slater--Condon algebra fixes $\IE(\mathrm{O})\approx\ERy$ to tens of
  per cent, with the $2p^4$ pairing penalty ($0.86\,$eV, to $6\%$ of
  the nitrogen--oxygen anomaly) supplying the drop onto the Rydberg
  scale.
\item The residual excess $\IE(\mathrm{O})/\ERy - 1 = 9.04\times10^{-4}$
  is $O(\aem)$ and coincidentally close to $3\aem/(8\pi)$, but that
  form is excluded as an additive condensate coupling shift by the
  hydrogen ionisation energy (Remark~\ref{P11:rem:oalpha_excess}); the
  excess is $2p^4$ correlation energy, condensate-downstream through
  the fixed Hamiltonian but not a coupling shift beyond it.
\end{itemize}
The parameter-free chain reaches the atomic scale but not the
coincidence's precision:
$(\TCMB,m_e)\xrightarrow{\text{Papers I--X}}(\aem,m_e)
\xrightarrow{\text{Cor.~\ref{P11:cor:rydberg}}}n(\ERy)=-12.102$,
with $\IE(\mathrm{O})\approx\ERy$ an atomic-structure fact.

\paragraph{Successive ionisation energies and the multi-layer spiral.}
Section~\ref{P11:sec:successive_IE} extends the first-IE analysis to
all successive ionisation energies.
The master formula (Proposition~\ref{P11:prop:master_formula}) places
IE$_k$ at tower level
$\nlvl{\IE_k}=\nlvl{\ERy}-\log_\vp Z_{\rm eff}(k)+\log_\vp n_{\rm qn}(k)$,
with no free parameters beyond $\aem$ and $m_e$.
Three shell completions lie near integer spiral levels: Pb~IE$_3$ at
$\Delta n=0.011$ from $-13$, Xe~IE$_3$ at $\Delta n=0.005$ from $-13$,
Xe~IE$_1$ at $\Delta n=0.0173$ from $-12$
(Proposition~\ref{P11:prop:shell_completion}).  A census of the
tabulated ionisation energies places such landings at or below the
chance rate, so these are individual placements and not a
noble-gas integer-marking rule.
The inner-shell pair O~IE$_6$--Na~IE$_5$ is near-degenerate at
$|\Delta n|=0.0021$ (Proposition~\ref{P11:prop:inner_degen}),
numerically comparable to the H--O first-IE result and arising from a
common effective charge $Z_{\rm eff}\approx6.37$ at equal principal
quantum number, though from different sub-shells.  Both this pair and
the near-integer shell-completion placements occur at rates consistent
with chance under their own trial counts, and are recorded as
placements rather than as evidence of spiral structure.  The spiral crossing predicted by this near-degeneracy
identifies the oxidation state at which O and Na become
Goldschmidt-substitutable.
Theorem~\ref{P11:thm:auc} establishes that spiral proximity $\Delta n$ predicts
Goldschmidt geochemical substitution compatibility with
$\mathrm{AUC}=0.857$ over all $\binom{118}{2}=6903$ element pairs ---
substantially better than ionic-radius proximity alone
($\mathrm{AUC}\approx0.72$).
At threshold $\Delta n<0.05$, $58\%$ of the 50 known pairs are recovered.
The code is archived at~\cite{Zenodo_P11}.

The framework does not predict individual ionisation energies
from first principles --- that requires solving the full
multi-electron Schr\"{o}dinger equation.
What it predicts is the \emph{scale} at which atomic physics lives:
the Rydberg is fixed, and all IEs are within a factor of
$\vp^2 \approx 2.6$ of it.
Chemistry is not derived from the condensate, but it is
\emph{calibrated} by the condensate through the two output
quantities $\aem$ and $m_e$.

The $\vp$-spiral proposed here is an alternative organisational
principle for the periodic table --- ordered by binding energy
rather than by electron configuration.
Its primary predictive application is crystal substitution:
Conjecture~\ref{P11:conj:spiral_pt} and Proposition~\ref{P11:prop:goldschmidt}
offer a falsifiable test against the Goldschmidt compatibility database.

%══════════════════════════════════════════════════════════════════════════════
\appendix
%══════════════════════════════════════════════════════════════════════════════

\section{Numerical Values}
\label{P11:app:numerics}

Key numerical values used throughout:
\begin{align}
  \vp &= (1+\sqrt{5})/2 = 1.6180339887\ldots \\
  \Lcond &= 1.18950\times10^{-4}\,\mathrm{eV} \\
  \aem^{-1} &= 137.03564 \quad (\text{framework; Paper~IV~}) \\
  m_e &= 0.51100\,\mathrm{MeV} \quad (\text{input}) \\
  \ERy &= \aem^2 m_e/2 = 13.6057\,\mathrm{eV} \\
  n(\ERy) &= -12.102 \\
  \Lcond\,\vp^{20} &= 1.7994\,\mathrm{eV} \\
  \vp^{20} &= L_{20} + \epsilon_{20} \approx 15127.000
    \quad (L_{20} = 15127 \text{ exact Lucas number}).
\end{align}
Lucas numbers $L_n = L_{n-1}+L_{n-2}$ with $L_0=2$, $L_1=1$:
$\ldots, L_{18}=9349$, $L_{20}=15127$, $L_{22}=39603$.
The deviation $\epsilon_{20} = \vp^{20} - L_{20} \approx -10^{-5}$
is exponentially small by the Binet formula.

\section{Interactive Three.js Visualisation}
\label{P11:app:code}

The $\vp$-spiral visualisation described in this paper is available
as a self-contained HTML/Three.js file.
It includes:
\begin{itemize}[noitemsep]
\item All particle masses, gauge bosons, and QCD scale
  on the 3D helix;
\item All 118 elements (90 naturally occurring + 28 synthetic)
  placed by first IE, coloured by periodic table group;
  toggle between natural-only and full-118 display;
\item A live ``levels per revolution'' slider ($N_\mathrm{arms}$)
  that redistributes all objects in real time,
  allowing exploration of different angular compressions;
\item Hover tooltips with exact values, tower levels, and
  framework context for every object;
\item Toggle controls for labels and element visibility.
\end{itemize}
The source is included as supplementary material.

%──────────────────────────────────────────────────────────────────────────────
\begin{thebibliography}{99}

\bibitem{Paper1}
F.~Manfredi,
\textit{The Density-Feedback Faddeev--Niemi Hopfion:
  Exact Algebraic Predictions for Standard-Model Parameters},
Zenodo (2026).
\href{https://doi.org/10.5281/zenodo.19342027}{doi:10.5281/zenodo.19342027}

\bibitem{Paper2}
F.~Manfredi,
\textit{BPS Structure and Fixed-Point Theorem for the
  Density-Feedback Faddeev--Niemi Hopfion},
Zenodo (2026).
\href{https://doi.org/10.5281/zenodo.19363491}{doi:10.5281/zenodo.19363491}

\bibitem{Paper3}
F.~Manfredi,
\textit{Two Constants from One Knot: The Fine Structure Constant and
  the Linking Scale as Exact Predictions of the Icosahedral WZW Condensate},
Zenodo (2026).
\href{https://doi.org/10.5281/zenodo.19478629}{doi:10.5281/zenodo.19478629}

\bibitem{Paper4}
F.~Manfredi,
\textit{Hopfion IV - The Hopf Spoke, WZW Fermion Mass Renormalization, and Three- and Four-Term Formulas for the Fine Structure Constant},
Preprint (2026).
\href{https://doi.org/10.5281/zenodo.19504729}{doi:10.5281/zenodo.19504729}

\bibitem{Paper5}
F.~Manfredi,
\textit{Colour Confinement, the QCD Scale, and Hadronic Structure
  from the Density-Feedback Hopfion},
Preprint (2026).
\href{https://doi.org/10.5281/zenodo.19638857}{doi:10.5281/zenodo.19638857}

\bibitem{Paper6}
F.~Manfredi,
\textit{The Golden-Ratio Tower: Fermion Scale Hierarchy from the Hopfion RG Fixed Point},
Preprint (2026).
\href{https://doi.org/10.5281/zenodo.19646945}{doi:10.5281/zenodo.19646945}

\bibitem{Paper7}
F.~Manfredi,
\textit{Emergent Gravity, Inflation, Dark Energy, and the Weak
  Equivalence Principle from the Density-Feedback Hopfion Condensate},
Preprint (2026).
\href{https://doi.org/10.5281/zenodo.19646953}{doi:10.5281/zenodo.19646953}

\bibitem{Paper8}
F.~Manfredi,
\textit{Geometry-Dependent Bell Violations from the Density-Feedback Hopfion},
Zenodo (2026).
\href{https://doi.org/10.5281/zenodo.18737070}{doi:10.5281/zenodo.18737070} 

\bibitem{Paper9}
F.~Manfredi,
\textit{Division Algebras, the Born Rule, and the Tsirelson Bound
  from the Density-Feedback Hopfion},
Zenodo (2026).
\href{https://doi.org/10.5281/zenodo.20075227}{doi:10.5281/zenodo.20075227}

\bibitem{Paper10}
F.~Manfredi,
\textit{Quantisation of the Hopfion Condensate: Hilbert Space, Liouville Measure, and the Born Rule},
Zenodo (2026).
\href{https://doi.org/10.5281/zenodo.20075279}{doi:10.5281/zenodo.20075279}

\bibitem{Paper12}
F.~Manfredi,
\textit{The Profile Normalisation Conjecture ---
  Proof via CMB Energy Identity and Two-Route Equivalence},
Zenodo (2026).
\href{https://doi.org/10.5281/zenodo.20210460}{doi:10.5281/zenodo.20210460}

\bibitem{Goldschmidt1937}
V.~M.~Goldschmidt,
\textit{The principles of distribution of chemical elements in
  minerals and rocks},
J.~Chem.\ Soc.\ \textbf{140}, 655--673 (1937).
\href{https://doi.org/10.1039/JR9370000655}{doi:10.1039/JR9370000655}

\bibitem{NIST_ASD}
A.~Kramida, Yu.~Ralchenko, J.~Reader, and NIST ASD Team,
\textit{NIST Atomic Spectra Database (version 5.11)},
National Institute of Standards and Technology, Gaithersburg, MD (2023).
\href{https://doi.org/10.18434/T4W30F}{doi:10.18434/T4W30F}

\bibitem{Slater1930}
J.~C.~Slater,
\textit{Atomic shielding constants},
Phys.\ Rev.\ \textbf{36}, 57--64 (1930).
\href{https://doi.org/10.1103/PhysRev.36.57}{doi:10.1103/PhysRev.36.57}

\bibitem{PDG2024}
R.~L.~Workman et al.\ (Particle Data Group),
\textit{Review of Particle Physics},
Prog.\ Theor.\ Exp.\ Phys.\ \textbf{2022}, 083C01 (2022);
updated 2024.
\href{https://doi.org/10.1093/ptep/ptac097}{doi:10.1093/ptep/ptac097}

\bibitem{Clementi1963}
E.~Clementi and D.~L.~Raimondi,
\textit{Atomic screening constants from SCF functions},
J.~Chem.\ Phys.\ \textbf{38}, 2686--2689 (1963).
\href{https://doi.org/10.1063/1.1733573}{doi:10.1063/1.1733573}

\bibitem{Burns1967}
R.~G.~Burns and W.~S.~Fyfe,
\textit{Trace element distribution rules and their significance},
Chemical Geology \textbf{2}, 89--104 (1967).
\href{https://doi.org/10.1016/0009-2541(67)90010-1}{doi:10.1016/0009-2541(67)90010-1}

\bibitem{Condon1935}
E.~U.~Condon and G.~H.~Shortley,
\textit{The Theory of Atomic Spectra},
Cambridge University Press (1935; reprinted 1951).
ISBN 9780521092098.

\bibitem{Zenodo_P11}
F.~Manfredi,
\textit{Code supplement for Paper~XI: $\vp$-spiral Goldschmidt ROC analysis
  and 118-element ionisation energy dataset},
Zenodo (2026).
\href{https://doi.org/10.5281/zenodo.20173763}{doi:10.5281/zenodo.20173763}

\bibitem{Bindi2009}
L.~Bindi, P.~J.~Steinhardt, N.~Yao, and P.~J.~Lu,
\textit{Natural Quasicrystals},
Science \textbf{324}, 1306 (2009).
\href{https://doi.org/10.1126/science.1170827}{doi:10.1126/science.1170827}

\end{thebibliography}

\end{document}
